\begin{theorem}[Euler number of a Kuhn cusp]\label{maps:src:F.1.1}
\StatementRecord{maps:src:F.1.1}{proved}{GPT-6 Astra (OpenAI, ChatGPT Pro mode), effort Pro, 2026-11-17; GPT-6 Astra (OpenAI, ChatGPT Pro mode), S6\_06\_BLIND, effort Pro, 2026-11-17; Claude Fable 5.2 (Anthropic), effort high, 16 November 2026; also checked in this writing by \CurrentChecker}{the hypotheses, notation and definitions stated here; the explicit constructions and calculations in the proof}


Let $g:M^8\to S^2$ be a proper triangulable map, a $T^6$-bundle
off finitely many values.  Suppose its unique non-quotient special
fibre is the rank-three Kuhn toroidal fibre with period-lattice
index $d$, and all its other special reduced fibres are free
finite quotients of $T^6$.  Then $\chi(M)=6d$.

\end{theorem}

\begin{proposition}[the diagonal obstruction]\label{maps:src:F.1.4}
\StatementRecord{maps:src:F.1.4}{proved}{GPT-6 Astra (OpenAI, ChatGPT Pro mode), effort Pro, 2026-11-17; Claude Fable 5.2 (Anthropic), effort high, 16 November 2026; also checked in this writing by \CurrentChecker}{the hypotheses, notation and definitions stated here; the explicit constructions and calculations in the proof}


If $A\in\operatorname{GL}_n(\mathbb Z)$ has no eigenvalue one,
every affine map $x\mapsto Ax+b$ of $\mathbb R^n/\mathbb Z^n$
has a fixed point.  Thus the diagonal elliptic finite monodromies
in Lead 7 admit no free affine lifts.

\end{proposition}

\begin{proposition}[the cusp normal form]\label{maps:src:F.2.1}
\StatementRecord{maps:src:F.2.1}{proved}{GPT-6 Astra (OpenAI, ChatGPT Pro mode), effort Pro, 2026-11-17; GPT-6 Astra (OpenAI, ChatGPT Pro mode), S6\_06\_BLIND, effort Pro, 2026-11-17; Claude Fable 5.2 (Anthropic), effort high, 16 November 2026; also checked in this writing by \CurrentChecker}{the hypotheses, notation and definitions stated here; the explicit constructions and calculations in the proof}


Under the integral index-one condition in F.2, the cusp matrix is
conjugate in $\operatorname{GL}_6(\mathbb Z)$ to
\[
 U=\begin{pmatrix}I_3&I_3\\0&I_3\end{pmatrix}.
\]
Its integral centralizer consists exactly of
$\left(\begin{smallmatrix}P&Q\\0&P\end{smallmatrix}\right)$
with $P\in\operatorname{GL}_3(\mathbb Z)$ and
$Q\in\operatorname{Mat}_{3\times3}(\mathbb Z)$.

\end{proposition}

\begin{proposition}[a complete finite-order filter]\label{maps:src:F.2.2}
\StatementRecord{maps:src:F.2.2}{proved}{GPT-6 Astra (OpenAI, ChatGPT Pro mode), S6\_06\_BLIND, effort Pro, 2026-11-17; Claude Fable 5.2 (Anthropic), effort high, 16 November 2026; also checked in this writing by \CurrentChecker}{the hypotheses, notation and definitions stated here; \cref{maps:src:F.1.4}}


The linear order of a finite matrix in $\operatorname{GL}_6(\mathbb Z)$
which admits a free affine lift belongs to
\[
 \{1,2,3,4,5,6,8,10,12\}.
 \tag{F.3}\label{maps:eq:F.3}
\]
There are 38 possible degree-six cyclotomic characteristic
polynomials with an eigenvalue one, and 19 of them commute
with some real complex structure.  These are necessary types,
not 38 or 19 conjugacy classes of admissible tuples.

\end{proposition}

\begin{proposition}[Counterexample: determination by peripheral data alone]\label{maps:src:F.2.3}
\StatementRecord{maps:src:F.2.3}{false as stated; correction proved}{Claude Fable 5.2 (Anthropic), effort high, 16 November 2026; also checked in this writing by \CurrentChecker}{the hypotheses, notation and definitions stated here; the explicit constructions and calculations in the proof}


Monodromy matrices and affine twist vectors alone do not
determine the homology of a compact glued manifold.  A cusp
filling, or an equivalent relative topological specification,
is additional data.

\end{proposition}

\begin{theorem}[the self-dual parity obstruction]\label{maps:src:F.3.1}
\StatementRecord{maps:src:F.3.1}{proved}{Claude Fable 5.2 (Anthropic), effort high, 16 November 2026; also checked in this writing by \CurrentChecker}{the hypotheses, notation and definitions stated here; \cref{maps:src:F.1.1}}


Suppose a compact smooth eight-manifold has the torus fibration
and the specified Kuhn cusp of \ref{maps:src:F.1.1} with index one, and its
other reduced fibres are free quotients of complex three-tori
by finite holomorphic affine actions.  If its rational
$H^1$ monodromy representation is self-dual, then it cannot
satisfy $b_1=b_2=0$.  In particular no symplectic-monodromy
datum of F.2 can give the three-connected fourfold required
in F.4, for either $k=2$ or $k=3$ (or any number $k$ of finite
points).

\end{theorem}

\begin{proposition}[the remaining rank patterns]\label{maps:src:F.3.2}
\StatementRecord{maps:src:F.3.2}{proved}{Claude Fable 5.2 (Anthropic), effort high, 16 November 2026; also checked in this writing by \CurrentChecker}{the hypotheses, notation and definitions stated here; the explicit constructions and calculations in the proof}


If a possibly non-self-dual holomorphic datum with the same
local models had $b_1=b_2=0$, put
$b=\dim H^0(B,\mathcal V^*)$.  For $k=2$, the only possible
unordered pairs $(r_1,r_2;b)$ are
\[
 (2,2;0),\qquad(2,4;2).
\]
For $k=3$, they are
\[
 (2,2,6;0),\ (2,4,4;0),\ (2,4,6;2),\ (4,4,4;2).
 \tag{F.4}\label{maps:eq:F.4}
\]

\end{proposition}

\begin{proposition}[resolution of the classification problem]\label{maps:src:F.2.4}
\StatementRecord{maps:src:F.2.4}{proved}{Claude Fable 5.2 (Anthropic), effort high, 16 November 2026; also checked in this writing by \CurrentChecker}{\cref{maps:src:F.2.16}, \cref{maps:src:F.2.30}, \cref{maps:src:F.2.31}, \cref{maps:src:F.2.35}, \cref{maps:src:F.2.36}, \cref{maps:src:F.2.37}, \cref{maps:src:F.2.38}, \cref{maps:src:F.2.39}}


The previously incomplete linear classification and all-twist
integral topology computation are completed by \ref{maps:src:F.2.31} and
\ref{maps:src:F.2.35}--39, with the fixed cusp filling and full marking
specified as required by \ref{maps:src:F.2.3}.  Every linear tuple has an
exact canonical representative and a returned unimodular
conjugator; the finite bound depends on the supplied input.
The unbounded trace family proves that no fixed global
height cutoff can replace this procedure.  All actual affine
orders, numerators, clutch vectors and cusp markings are
included in the general fundamental-group and integral
homology algorithms.  The complete constructive survivor
list, its integral tests and the common oriented smooth
type are given in \ref{maps:src:F.2.30}.

This conclusion does not claim an executed enumeration of
the infinite general class set.  The exhaustive finite
involution catalogue, the general comparison tests, the
three full topology integrations and the unified wrapper
checks have the exact executed scopes stated in \ref{maps:src:F.2.35}--39.
The earlier obstruction record is retained verbatim in
\path{records/fourfolds_F2_4_superseded.tex}; its finite-order
filter and homology-level obstructions were superseded by
the complete canonical and chain constructions.


\end{proposition}

\begin{theorem}[a non-self-dual topological candidate]\label{maps:src:F.2.5}
\StatementRecord{maps:src:F.2.5}{proved}{GPT-6 Astra (OpenAI, ChatGPT Pro mode), effort Pro, 2026-11-17; Claude Fable 5.2 (Anthropic), effort high, 2026-11-16; GPT-6 Astra (OpenAI, ChatGPT Pro mode), S6\_06\_BLIND, effort Pro, 2026-11-17; also checked in this writing by \CurrentChecker}{the hypotheses, notation and definitions stated here; \cref{maps:src:F.2.1}; \ExternalPaperRef{theorem~}{6.3}{III}; \S\ref{sec:V-known-inputs}}


There is a closed smooth oriented eight-manifold with the prescribed
index-one Kuhn cusp, two free affine fillings of multiplicities $5,4$,
trivial fundamental group, and rational Betti numbers
\[
 (1,0,0,0,4,0,0,0,1). \tag{F.5}\label{maps:eq:F.5}
\]
It admits an almost complex structure.  Its signature is two and
its middle intersection lattice modulo torsion is
$3\langle1\rangle\oplus\langle-1\rangle$.
Integral $H_2,H_3$ and a holomorphic period realization are not
asserted here.

\end{theorem}

\begin{theorem}\label{night:P1.1}
\StatementRecord{night:P1.1}{proved}{S6\_06\_BLIND: GPT-6 Astra (OpenAI, ChatGPT Pro mode, effort Pro), recorded re-derivation in the checkpoint dated 18 November 2026; also checked in this writing by \CurrentChecker}{Definitions and the displayed constructions}
The finite homology monodromies have orders $5,4$, respectively, and
characteristic polynomials
$(T-1)^2\Phi _5(T)$ and $(T-1)^2(T+1)^2(T^2+1)$.
For $N=M_0-I$ we have $N^2=0$, $\operatorname{rank}N=3$, and
nonzero Smith factors $1,1,1$.  Moreover
\[
 \Lambda^{\langle M_1,M_2\rangle}=0,\qquad
 \Lambda_{\langle M_1,M_2\rangle}=\mathbb Z,
\]
where the quotient homomorphism is $\ell(x)=x_0$.
On the cohomology lattice the corresponding groups are $\mathbb Ze_0$
and $\mathbb Z/2$.  Every nonidentity power of each specified affine
action $\tau_i$ is free.
\end{theorem}

\begin{theorem}\label{night:P1.3}
\StatementRecord{night:P1.3}{proved}{S6\_06\_BLIND: GPT-6 Astra (OpenAI, ChatGPT Pro mode, effort Pro), recorded re-derivation in the checkpoint dated 18 November 2026; also checked in this writing by \CurrentChecker}{Definitions and the displayed constructions}
The two finite filling pieces have first integral homology
$\mathbb Z^2$ and $\mathbb Z^2\oplus\mathbb Z/2$.
\end{theorem}

\begin{lemma}[The actual order-four torsion transgression]
\label{mapsaudit:order-four-torsion-transgression}
\StatementRecord{mapsaudit:order-four-torsion-transgression}{proved}{GPT-6 Astra (OpenAI, ChatGPT Pro mode, effort Pro), S6\_09 writing, 17 November 2026; also checked in this writing by \CurrentChecker}{\cref{maps:src:F.2.5}; the displayed cellular and boundary maps}

Use the order-four filling and the marked matrices of \cref{maps:src:F.2.5}.
Let $T=\bigwedge^{q-1}D^t$ on the integral vertical cochains of the
five-torus. If $w$ represents a torsion element of
$\operatorname{coker}(T-I)$, its point-supported degree-$q$ class has
Leray differential equal, up to the one boundary-orientation sign, to
the image of $w$ in the global degree-$(q-1)$ coinvariants. There is
no multiplicity-four factor. For each of $q=2,3,4$, that image is
the nonzero class in the global $\mathbb Z/2$ quotient.
\end{lemma}

\begin{theorem}[integral three-connectivity of the candidate]\label{maps:src:F.2.6}
\StatementRecord{maps:src:F.2.6}{proved}{GPT-6 Astra (OpenAI, ChatGPT Pro mode), effort Pro, 2026-11-17; Claude Fable 5.2 (Anthropic), effort high, 2026-11-17; Claude Fable 5.2 (Anthropic), effort high, 2026-11-16; GPT-6 Astra (OpenAI, ChatGPT Pro mode), S6\_06\_BLIND, effort Pro, 2026-11-17; also checked in this writing by \CurrentChecker}{\cref{mapsaudit:order-four-torsion-transgression}; the hypotheses, notation and definitions stated here; \cref{maps:src:F.2.5}; \cref{maps:eq:F.5}; \ExternalPaperRef{theorem~}{6.3}{III}; \cref{maps:eq:F.6}}


The manifold of Theorem \ref{maps:src:F.2.5} has free integral homology with
Betti numbers \eqref{maps:eq:F.5}.  In particular $H_2=H_3=0$, and it is
three-connected.

\end{theorem}

\begin{theorem}\label{night:P1.10}
\StatementRecord{night:P1.10}{proved}{S6\_06\_BLIND: GPT-6 Astra (OpenAI, ChatGPT Pro mode, effort Pro), recorded re-derivation in the checkpoint dated 18 November 2026; also checked in this writing by \CurrentChecker}{\cref{night:P1.11}, \cref{night:P1.9}}
For every real parameter $c>-h_{\min}$ the corrected construction gives
a compact smooth complex fourfold $M_c$ and a proper holomorphic map
to $\mathbb P^1$ with the two specified multiple fibres and the
index-one Kuhn cusp.  With its complex orientation,
$\pi_1(M_c)=1$ and $\chi(M_c)=6$.
\end{theorem}

\begin{theorem}\label{night:P1.11}
\StatementRecord{night:P1.11}{proved}{S6\_06\_BLIND: GPT-6 Astra (OpenAI, ChatGPT Pro mode, effort Pro), recorded re-derivation in the checkpoint dated 18 November 2026; also checked in this writing by \CurrentChecker}{\cref{night:P1.9}}
The index-one central cusp fibre $C$ has torsion-free integral homology
with ranks $(1,3,9,9,10,3,1)$.
The collapse $\kappa:T^6\to C$ has primitive image in every degree;
it is onto except in degree four, where its cokernel is $\mathbb Z$.
Its image ranks are $(1,3,9,9,9,3,1)$.
In every degree its kernel is the saturation of
$(\bigwedge^kM_0-I)\bigwedge^k\Lambda$.
\end{theorem}

\begin{theorem}\label{night:P1.12}
\StatementRecord{night:P1.12}{proved}{S6\_06\_BLIND: GPT-6 Astra (OpenAI, ChatGPT Pro mode, effort Pro), recorded re-derivation in the checkpoint dated 18 November 2026; also checked in this writing by \CurrentChecker}{\cref{night:P1.11}}
The fourfold $M_c$ has torsion-free integral homology
\[
H_k(M_c;\mathbb Z)=
\begin{cases}
\mathbb Z,&k=0,8,\\
\mathbb Z^4,&k=4,\\
0,&k=1,2,3,5,6,7.
\end{cases}
\]
In particular it is three-connected.
\end{theorem}

\begin{theorem}\label{night:P1.19}
\StatementRecord{night:P1.19}{proved}{S6\_06\_BLIND: GPT-6 Astra (OpenAI, ChatGPT Pro mode, effort Pro), recorded re-derivation in the checkpoint dated 18 November 2026; also checked in this writing by \CurrentChecker}{\cref{night:P1.11}, \cref{night:P1.12}, \cref{night:P1.9}}
There is an independent computation from the full global cellular
complex.  Its chain ranks in degrees zero through eight are
\[
 (4,26,87,162,200,156,79,23,3),
\]
and its differentials, after the explicit integral chain equivalences
below, are the matrices of case one in
\path{code/part1/global_ext_cases.json}.
Their integral Smith calculation gives
\[
 H_k(M;\mathbb Z)=
 \begin{cases}
 \mathbb Z,&k=0,8,\\
 \mathbb Z^4,&k=4,\\
 0,&\text{otherwise}.
 \end{cases}
\]
This calculation uses local Poincar\'e--Lefschetz duality of the
smooth cusp cap to determine its secondary attaching class.  It does
not use global Poincar\'e duality or the homology conclusion of P1.12.
In particular, the map obtained by setting the whole cusp suspension
chain term to zero is not the actual relative cusp model; its precise
missing term is accounted for here.
\end{theorem}

\begin{theorem}\label{night:P1.4}
\StatementRecord{night:P1.4}{proved}{S6\_06\_BLIND: GPT-6 Astra (OpenAI, ChatGPT Pro mode, effort Pro), recorded re-derivation in the checkpoint dated 18 November 2026; also checked in this writing by \CurrentChecker}{Definitions and the displayed constructions}
The genus-two pencil has one nodal fibre away from $0,\infty$, at
$\lambda_*=-256/3125$.  Good reduction at $0,\infty$ requires root
orders $8,5$.  On this root orbifold its Hodge bundle is
$\mathcal H=T^3\oplus T$, where
\[
T=\mathcal O\bigl((2/5)[\infty]-(3/8)[0]\bigr),\qquad
\deg T=1/40,
\]
and $T^3$ is the line generated meromorphically by $dx/y$.
\end{theorem}

\begin{theorem}\label{night:P1.5}
\StatementRecord{night:P1.5}{false as stated; correction proved}{S6\_06\_BLIND: GPT-6 Astra (OpenAI, ChatGPT Pro mode, effort Pro), recorded re-derivation in the checkpoint dated 18 November 2026; also checked in this writing by \CurrentChecker}{Definitions and the displayed constructions}
The formula in the prompt with $P$ the image of
$\ell_0\wedge\mathcal H$, and $\mathcal H$ the rank-two Hodge bundle,
cannot produce a rank-three subbundle.  The required rank-three image
is instead the image of $\ell_0\wedge\mathcal H_{\mathrm{dR}}$, where
$\mathcal H_{\mathrm{dR}}$ is the full rank-four de Rham bundle.
\end{theorem}

\begin{theorem}[the geometric rank-five quotient]\label{maps:src:F.3.5}
\StatementRecord{maps:src:F.3.5}{proved}{GPT-6 Astra (OpenAI, ChatGPT Pro mode), effort Pro, 2026-11-17; GPT-6 Astra (OpenAI, ChatGPT Pro mode), S6\_06\_BLIND, effort Pro, 2026-11-17; Claude Fable 5.2 (Anthropic), effort high, 2026-11-16; also checked in this writing by \CurrentChecker}{the hypotheses, notation and definitions stated here; \cref{maps:src:F.2.5}}


Put $\lambda_*=-256/3125$ and
$U=\mathbb P^1\setminus\{0,\infty,\lambda_*\}$.  The smooth projective
models of
\[
 C_\lambda:\quad y^2=x^5+\lambda^3x+\lambda^4
 \tag{F.7}\label{maps:eq:F.7}
\]
give a polarized rank-four rational cohomology system $\mathcal H$ on
$U$.  Its primitive exterior square is rationally isomorphic to
the rank-five quotient of the system with matrices $A,B$ in
Theorem \ref{maps:src:F.2.5} by its invariant first coordinate line.
On the orbifold $\mathcal B'$ obtained by taking roots of orders
five at infinity and eight at zero, the extended Hodge bundle is
\[
 \mathcal E=\mathcal T^3\oplus\mathcal T,\qquad
 \mathcal T=\mathcal O_{\mathcal B'}(2D_\infty-3D_0),\qquad
 \deg\mathcal T=\frac1{40}.
 \tag{F.8}\label{maps:eq:F.8}
\]
Here $D_\infty,D_0$ are the reduced root divisors; the extension at
$\lambda_*$ is the nilpotent canonical extension.

\end{theorem}

\begin{theorem}\label{night:P1.2}
\StatementRecord{night:P1.2}{proved}{S6\_06\_BLIND: GPT-6 Astra (OpenAI, ChatGPT Pro mode, effort Pro), recorded re-derivation in the checkpoint dated 18 November 2026; also checked in this writing by \CurrentChecker}{Definitions and the displayed constructions}
The lower blocks of $A,B$ are precisely the quotient representations of
$\bigwedge^2(C_f^{-1}),\bigwedge^2C_g$ in the supplied basis $Q$.
The integral quotient is the dual of the primitive exterior-square
lattice.  The natural map from the primitive lattice into that quotient
has cokernel $\mathbb Z/2$; these two integral representations are not
isomorphic.  Their rationalizations are isomorphic by that same map.
\end{theorem}

\begin{theorem}\label{night:P1.7}
\StatementRecord{night:P1.7}{proved}{S6\_06\_BLIND: GPT-6 Astra (OpenAI, ChatGPT Pro mode, effort Pro), recorded re-derivation in the checkpoint dated 18 November 2026; also checked in this writing by \CurrentChecker}{\cref{night:P1.2}}
The integral $H^1$ monodromy of the pencil has a companion marking with
finite pair $(C_f^{-1},C_g)$, and cusp product $C_fC_g^{-1}$.
Its exterior-square quotient is consequently the integral quotient
representation appearing in P1.2.
\end{theorem}

\begin{theorem}[a global nonpolarized period family]\label{maps:src:F.3.6}
\StatementRecord{maps:src:F.3.6}{proved}{GPT-6 Astra (OpenAI, ChatGPT Pro mode), effort Pro, 2026-11-17; GPT-6 Astra (OpenAI, ChatGPT Pro mode), S6\_06\_BLIND, effort Pro, 2026-11-17; Claude Fable 5.2 (Anthropic), effort high, 2026-11-16; also checked in this writing by \CurrentChecker}{the hypotheses, notation and definitions stated here; \cref{maps:eq:F.11}; \cref{maps:eq:F.14}; \cref{maps:eq:F.16}; \cref{maps:eq:F.19}; \cref{maps:eq:F.12}; \cref{maps:src:F.2.1}; \cref{maps:eq:F.22}; \cref{maps:eq:F.20}; \cref{maps:eq:F.21}}


The integral rank-six system with cohomology monodromies $A,B$
admits a holomorphic subbundle $\mathcal F$ of rank three over
$U$ satisfying $\mathcal F\cap\overline{\mathcal F}=0$.
It extends holomorphically on the finite root charts.  At the
cusp an integral marking puts its periods in the form
\[
 Z(t)=\frac{\log t}{2\pi i}I_3+C(t),
 \tag{F.13}\label{maps:eq:F.13}
\]
where $C$ is holomorphic at zero and bounded there.

\end{theorem}

\begin{theorem}\label{night:P1.8}
\StatementRecord{night:P1.8}{proved}{S6\_06\_BLIND: GPT-6 Astra (OpenAI, ChatGPT Pro mode, effort Pro), recorded re-derivation in the checkpoint dated 18 November 2026; also checked in this writing by \CurrentChecker}{\cref{night:P1.5}}
With the corrected rank-three bundle $P$ of P1.5, the supplied rank-six
system admits a holomorphic lifted family
$F_c=(\sigma+ic\alpha e)(P)$ on the punctured base.
For a specified splitting $\sigma$ there is a real analytic function
$h$ with finite attained minimum $h_{\min}$, and the exact domain of
real parameters giving $F_c\cap\overline{F_c}=0$ everywhere is
$c>-h_{\min}$.
At the cusp its periods have the form
$Z(q)=(\log q/(2\pi i))I_3+C(q)$ with $C$ holomorphic and bounded.
The construction extends to the two finite root charts and descends to
orders $5,4$ on the rank-six system.
\end{theorem}

\begin{theorem}\label{night:P1.9}
\StatementRecord{night:P1.9}{proved}{S6\_06\_BLIND: GPT-6 Astra (OpenAI, ChatGPT Pro mode, effort Pro), recorded re-derivation in the checkpoint dated 18 November 2026; also checked in this writing by \CurrentChecker}{Definitions and the displayed constructions}
The periodic Kuhn fan, in the marking $S$, gives a smooth proper
holomorphic filling of the periods $Z(q)=sI_3+C(q)$ for a sufficiently
small disk.  Its central fibre is reduced with normal crossings.
Modulo the index-one period lattice it has respectively $1,7,12,6$
strata of dimensions $3,2,1,0$, each isomorphic to the corresponding
complex torus $(\mathbb C^*)^d$, and Euler number $6$.
An index-$d$ period sublattice gives $d$ times these stratum counts and
Euler number $6d$.
\end{theorem}

\begin{theorem}[compact holomorphic filling]\label{maps:src:F.3.7}
\StatementRecord{maps:src:F.3.7}{proved}{GPT-6 Astra (OpenAI, ChatGPT Pro mode), effort Pro, 2026-11-17; GPT-6 Astra (OpenAI, ChatGPT Pro mode), S6\_06\_BLIND, effort Pro, 2026-11-17; Claude Fable 5.2 (Anthropic), effort high, 16 November 2026; also checked in this writing by \CurrentChecker}{the hypotheses, notation and definitions stated here; \cref{maps:src:F.2.6}; \cref{maps:src:F.2.5}; \cref{maps:eq:F.24}; \cref{maps:eq:F.13}; \ExternalPaperRef{equation~}{{E.4}}{III}; \ExternalPaperRef{equation~}{{E.5}}{III}; \ExternalPaperRef{theorem~}{6.2}{III}; \ExternalPaperRef{theorem~}{6.3}{III}}


The smooth eight-manifold of Theorem \ref{maps:src:F.2.6} carries an integrable
complex structure.  There is a proper holomorphic map to
$\mathbb P^1$ whose regular fibres are complex three-tori,
whose two finite reduced fibres are the free affine quotients
of orders five and four specified in Theorem \ref{maps:src:F.2.5}, and whose
remaining fibre is the index-one Kuhn degeneration.

\end{theorem}

\begin{theorem}\label{night:P1.13}
\StatementRecord{night:P1.13}{proved}{S6\_06\_BLIND: GPT-6 Astra (OpenAI, ChatGPT Pro mode, effort Pro), recorded re-derivation in the checkpoint dated 18 November 2026; also checked in this writing by \CurrentChecker}{\cref{night:P1.12}; \S\ref{sec:V-known-inputs}; \cite[equations (1.2)--(1.4), pp.~3--4]{ffEsoleKang}; \cite[Theorem 2.2.1, p.~61]{ffPetersSterk}}
With the complex orientation the fourfold has
\[
\sigma(M_c)=2,\qquad H^4(M_c;\mathbb Z)\cong I_{3,1},\qquad
w_2(M_c)=0,
\]
and there is a diagonal integral basis in which
$p_1/2=(1,1,1,1)$.
Its Chern numbers are
$c_4=6$, $c_2^2=2$, and
$c_1^4=c_1^2c_2=c_1c_3=0$; its Pontryagin numbers are
$p_1^2=8$, $p_2=14$.
\end{theorem}

\begin{theorem}[a complex structure on the stated connected sum]\label{maps:src:F.4.1}
\StatementRecord{maps:src:F.4.1}{proved}{GPT-6 Astra (OpenAI, ChatGPT Pro mode), effort Pro, 2026-11-17; Claude Fable 5.2 (Anthropic), effort high, 2026-11-17; Claude Fable 5.2 (Anthropic), effort high, 2026-11-15; GPT-6 Astra (OpenAI, ChatGPT Pro mode), S6\_06\_BLIND, effort Pro, 2026-11-17; Claude Fable 5.2 (Anthropic), effort high, 2026-11-16; also checked in this writing by \CurrentChecker}{the hypotheses, notation and definitions stated here; \cref{maps:src:F.3.7}; \cref{maps:src:F.2.6}; \S\ref{sec:V-known-inputs}; the explicit integral reflections below; \cite[Corollary 1.4, p.~4]{ffMandal}; \cite[Definition 3.6 and Theorem 3.7(1), pp.~12--14]{IVVWallForm}; \cite[Theorem A, p.~178]{IVVKramerStolz}}


The compact complex fourfold of Theorem \ref{maps:src:F.3.7} is three-connected,
has Euler number six, and is orientation-preservingly diffeomorphic to
\[
 3\mathbb H\mathbb P^2\ \#\ (-\mathbb H\mathbb P^2).
 \tag{F.25}\label{maps:eq:F.25}
\]
Its integral middle lattice is $I_{3,1}$ and, in a diagonal integral
basis, $p_1/2=(1,1,1,1)$.  In particular the conclusion is an exact
smooth diffeomorphism, with no remaining homotopy-sphere ambiguity.

\end{theorem}

\begin{theorem}\label{night:P1.15}
\StatementRecord{night:P1.15}{proved}{S6\_06\_BLIND: GPT-6 Astra (OpenAI, ChatGPT Pro mode, effort Pro), recorded re-derivation in the checkpoint dated 18 November 2026; also checked in this writing by \CurrentChecker}{\cref{night:P1.1}, \cref{night:P1.8}}
The rank-six flat extension, as a complex flat bundle with its invariant
line, has the following explicit algebraic connection model.  Set
$\omega_k=x^kdx/y$, $0\leq k\leq3$, and use the quotient frame
\[
(w_2,w_1,w_0,w_{-1},w_{-2})
 =(\omega_{01},\omega_{02},\omega_{03},\omega_{13},\omega_{23}),
\qquad \omega_{12}=-3\omega_{03}.
\]
These are meromorphic frames of
$(L^2,L,\mathcal O,L^{-1},L^{-2})$, where $L=T^2$ on the root
orbifold of orders $5,4$.  If a frame is written as a row of vectors,
its derivative is right multiplication by
\[
\Gamma_6=\begin{pmatrix}0&b\\0&\Gamma_W\end{pmatrix},\qquad
b=\frac1{\lambda-\lambda_*}
 \left(-\frac{25}{6\lambda^2},-\frac{125}{48\lambda},0,1,0\right),
\tag{1.15.1}
\]
where
\[
\resizebox{\textwidth}{!}{$\displaystyle
\Gamma_W=\frac1{3125\lambda+256}
\begin{pmatrix}
-8(625\lambda+48)/\lambda&-8&10\lambda&75\lambda^2/2&0\\
-125/(2\lambda)&-4(625\lambda+48)/\lambda&8&0&75\lambda^2/2\\
300/\lambda^2&-375/\lambda&0&-48&60\lambda\\
120/\lambda^3&0&125/(2\lambda)&4(625\lambda+48)/\lambda&-8\\
0&120/\lambda^3&50/\lambda^2&-125/(2\lambda)&8(625\lambda+48)/\lambda
\end{pmatrix}.$}
\tag{1.15.2}
\]
This model specifies a canonical algebraic splitting, denoted
$\sigma_b$.  The compatible contraction used for the real parameter is
$\alpha=(625\pi i/128)\alpha_{\rm alg}$, where
$\alpha_{\rm alg}(w_0)=4/3$ and it vanishes on $w_1,w_2$.
\end{theorem}

\begin{theorem}\label{night:P1.16}
\StatementRecord{night:P1.16}{proved}{S6\_06\_BLIND: GPT-6 Astra (OpenAI, ChatGPT Pro mode, effort Pro), recorded re-derivation in the checkpoint dated 18 November 2026; also checked in this writing by \CurrentChecker}{\cref{night:P1.15}, \cref{night:P1.18}, \cref{night:P1.8}}
In the canonical splitting $\sigma_b$ of P1.15,
\[
\begin{aligned}
-1.0342766224163217994942488453783461&<h_b(1),\\
h_b(1)&<-1.0342766224163217994942488453783460.
\end{aligned}
\]
In particular $h_b(1)<-1$.  The parameter $c=5/(3\pi)$ is outside the
global transversality domain.  The lifted period matrices at
$\lambda=1,2$ have been evaluated by two different continuation methods
agreeing to more than thirty decimal digits.
\end{theorem}

\begin{theorem}\label{night:P1.18}
\StatementRecord{night:P1.18}{proved}{S6\_06\_BLIND: GPT-6 Astra (OpenAI, ChatGPT Pro mode, effort Pro), recorded re-derivation in the checkpoint dated 18 November 2026; also checked in this writing by \CurrentChecker}{\cref{night:P1.15}, \cref{night:P1.16}, \cref{night:P1.8}}
In the algebraic splitting $\sigma_b$, the function $h_b$ has its
unique global minimum at the finite-root point $\lambda=\infty$.
The minimum is the exact Frobenius period constant
\[
 h_{b,\min}=-\frac{96}{625\pi}\operatorname{Re}(XH_\infty)_{03},
 \qquad
 h_{b,\min}=-1.0945402503114027801162682630549457037\ldots,
\tag{1.18.1}
\]
where $X$ is the normalized integral intertwiner of P1.15 and
$H_\infty$ is the regular infinity-root limit of its fundamental
solution, normalized to be the identity at $\lambda=1$ before
multiplication by $X$.  Formula (1.18.1), the rational connection,
and those two normalizations specify the constant exactly.
A certified enclosure has radius $4.04\cdot10^{-56}$ about the
56-digit midpoint retained in
\path{code/part1/minimum_certificate_summary.json}.
Consequently the exact domain of completed transverse periods is
\[
 c>-h_{b,\min},\qquad
 \operatorname{Im}\xi>-h_{b,\min}
\]
for the real and complex parameters, respectively.
The member $c=10$ used in the numerical period checks is therefore
admissible on the whole completed base.

For the point-normalized splitting of P1.8, continued from
$1$ to $\lambda_*/2$ by the polygon
$1,1+i,\lambda_*/2+i,\lambda_*/2$, the exact comparison constant
$a_{\rm split}$ of (1.15.4) has certified value
\[
\begin{split}
 \operatorname{Re}a_{\rm split}
   &=-0.7099885679126427771723473455012308149\ldots,\\
 \operatorname{Im}a_{\rm split}
   &=-0.9134017059364153362143229223828373269\ldots.
\end{split}
\]
Thus the minimum in that splitting is
$-0.1811385443749874439019453406721083768\ldots$,
attained at the same base point.  Its real parameter interval is
\[
 c>0.1811385443749874439019453406721083768\ldots.
\]
The full parameter comparison remains
$\xi_b=\xi_{\rm original}-a_{\rm split}$; both its real and imaginary
parts are necessary to identify the two normalized families.
\end{theorem}

\begin{theorem}\label{night:P1A.1}
\StatementRecord{night:P1A.1}{proved}{S6\_06\_BLIND: GPT-6 Astra (OpenAI, ChatGPT Pro mode, effort Pro), recorded re-derivation in the checkpoint dated 18 November 2026; also checked in this writing by \CurrentChecker}{\cref{night:P1.10}, \cref{night:P1.8}, \cref{night:P1.9}}
For every admissible member, its vertical translation bundle on
$\widehat B$ and its first coherent direct image are
\[
 E_0=F^\vee=\mathcal O\oplus L^{-1}\oplus L^{-2},\qquad
 G=R^1h_*\mathcal O_M=\mathcal O\oplus L^{-1}\oplus L^{-2}.
\]
Moreover $R^qh_*\mathcal O_M=\bigwedge^qG$ for $0\le q\le3$,
$K_M=h^*L^{-6}$, and
\[
 (h^0(\mathcal O_M),\ldots,h^4(\mathcal O_M))=(1,1,0,0,0).
\]
\end{theorem}

\begin{theorem}\label{night:P1A.2}
\StatementRecord{night:P1A.2}{proved}{S6\_06\_BLIND: GPT-6 Astra (OpenAI, ChatGPT Pro mode, effort Pro), recorded re-derivation in the checkpoint dated 18 November 2026; also checked in this writing by \CurrentChecker}{\cref{night:P1.15}, \cref{night:P1A.1}}
The complete Hodge table of every admissible fourfold is the
following, with rows indexed by $p$ and columns by $q$:
\[
\begin{array}{c|ccccc}
h^{p,q}&0&1&2&3&4\\\hline
0&1&1&0&0&0\\
1&0&2&1&0&0\\
2&0&1&6&1&0\\
3&0&0&1&2&0\\
4&0&0&0&1&1
\end{array}
\]
\end{theorem}

\begin{theorem}\label{night:P1A.3}
\StatementRecord{night:P1A.3}{proved}{S6\_06\_BLIND: GPT-6 Astra (OpenAI, ChatGPT Pro mode, effort Pro), recorded re-derivation in the checkpoint dated 18 November 2026; also checked in this writing by \CurrentChecker}{Definitions and the displayed constructions}
For every admissible fourfold,
\[
 (h^0(T_M),h^1(T_M),h^2(T_M),h^3(T_M),h^4(T_M))
 =(1,1,1,2,0).
\]
The complex parameter family has a nonzero Kodaira--Spencer map
at every member and supplies a smooth one-dimensional miniversal
Kuranishi germ.
\end{theorem}

\begin{theorem}\label{night:P1A.4}
\StatementRecord{night:P1A.4}{proved}{S6\_06\_BLIND: GPT-6 Astra (OpenAI, ChatGPT Pro mode, effort Pro), recorded re-derivation in the checkpoint dated 18 November 2026; also checked in this writing by \CurrentChecker}{\cref{night:P1.1}}
The fourfold has algebraic dimension one and
\[
\Aut^0(M)=\mathbb C,\qquad
\Aut(M)=\mathbb C\times\mathbb Z/2.
\]
In a fixed holomorphic splitting, two members of the full
admissible complex parameter half-plane are biholomorphic
exactly when their parameters are equal.
\end{theorem}

\begin{theorem}\label{night:P1.17}
\StatementRecord{night:P1.17}{proved}{S6\_06\_BLIND: GPT-6 Astra (OpenAI, ChatGPT Pro mode, effort Pro), recorded re-derivation in the checkpoint dated 18 November 2026; also checked in this writing by \CurrentChecker}{\cref{night:P1.1}}
For the cusp marking $S$ fixed in P1.1 and the corresponding Kuhn fan,
none of the completed fourfolds in this family admits an
anti-holomorphic automorphism.  In particular none admits a real
structure.  More strongly, two members with this same cusp marking
are never anti-biholomorphic.
\end{theorem}

\begin{theorem}\label{night:P1A.5}
\StatementRecord{night:P1A.5}{proved}{S6\_06\_BLIND: GPT-6 Astra (OpenAI, ChatGPT Pro mode, effort Pro), recorded re-derivation in the checkpoint dated 18 November 2026; also checked in this writing by \CurrentChecker}{\ExternalPaperRef{theorem~}{17.4}{I}}
There is no such analytic extension across any finite boundary
point of $\mathcal H$, or across infinity. The same conclusion
holds after any finite ramified change of the local boundary
coordinate. In particular neither end of the real interval
$c_*<c<+\infty$ has a proper Hausdorff analytic extension in
this sense.
\end{theorem}

\begin{theorem}\label{night:P1A.6}
\StatementRecord{night:P1A.6}{proved}{S6\_06\_BLIND: GPT-6 Astra (OpenAI, ChatGPT Pro mode, effort Pro), recorded re-derivation in the checkpoint dated 18 November 2026; also checked in this writing by \CurrentChecker}{\cref{night:P1A.3}, \cref{night:P1A.4}}
A holomorphic punctured-disc family whose fibres are all
members of this fourfold family has an interior parameter
limit. After a base change of degree at most two, it is
isomorphic to the pullback of the constructed family by
a holomorphic map of the full disc to $\mathcal H$, and
therefore has a proper smooth filling. Such a family
cannot approach either parameter end.
\end{theorem}

\begin{theorem}[the orbifold quotient]\label{maps:src:F.2.7}
\StatementRecord{maps:src:F.2.7}{proved}{Claude Fable 5.2 (Anthropic), effort high, 2026-11-17; Claude Fable 5.2 (Anthropic), effort high, 16 November 2026; also checked in this writing by \CurrentChecker}{the hypotheses, notation and definitions stated here; the explicit constructions and calculations in the proof}


For the specified torus fillings in F.2, with affine
multiplicities $m_1,\ldots,m_k$, the fundamental group of the
total space has a canonical surjection onto
\[
 \langle x_1,\ldots,x_k\mid
       x_i^{m_i}=1,\ x_1\cdots x_k=1\rangle.
 \tag{F.27}\label{maps:eq:F.27}
\]
Consequently a simply connected example with $k=2$ requires
$\gcd(m_1,m_2)=1$.  There is no simply connected example with
$k=3$ and $m_1,m_2,m_3\geq2$.  An affine multiplicity-one
point is an ordinary smooth point and reduces to the problem
with one fewer multiple fibre.

\end{theorem}

\begin{proposition}[exclusion of an involutory finite part]\label{maps:src:F.2.8}
\StatementRecord{maps:src:F.2.8}{proved}{Claude Fable 5.2 (Anthropic), effort high, 16 November 2026; also checked in this writing by \CurrentChecker}{the hypotheses, notation and definitions stated here; \cref{maps:src:F.2.7}; \cref{maps:src:F.2.15}; \cref{maps:src:F.3.1}; \cref{maps:eq:F.3}}


A simply connected datum with $b_2=0$, $k=2$ and the stated
finite affine fillings has neither finite linear
monodromy of order one or two.  Its unordered pair of finite
linear orders must therefore belong to
\[
 \{(3,4),(3,5),(3,8),(3,10),(4,5),
       (5,6),(5,8),(5,12)\}.
 \tag{F.28}\label{maps:eq:F.28}
\]
The affine multiplicities may still be larger than these
linear orders.

\end{proposition}

\begin{proposition}[the only first-cohomology pattern for two points]\label{maps:src:F.2.9}
\StatementRecord{maps:src:F.2.9}{proved}{Claude Fable 5.2 (Anthropic), effort high, 16 November 2026; also checked in this writing by \CurrentChecker}{the hypotheses, notation and definitions stated here; \cref{maps:src:F.3.1}; \cref{maps:src:F.2.15}}


For a simply connected two-multiple-fibre datum with $b_2=0$,
the rational
cohomology system $V$ has one invariant, no coinvariants,
and fixed rank two at each finite generator.  Its quotient
$W=V/\mathbb Qe$ by the invariant line has neither invariants
nor coinvariants.  If one finite characteristic polynomial
on $W$ is $(T-1)\Phi_d(T)$ with $\deg\Phi_d=4$, then $W$
is irreducible over $\mathbb Q$.

\end{proposition}

\begin{theorem}[exclusion of three further order pairs]\label{maps:src:F.2.10}
\StatementRecord{maps:src:F.2.10}{proved}{Claude Fable 5.2 (Anthropic), effort high, 16 November 2026; also checked in this writing by \CurrentChecker}{the hypotheses, notation and definitions stated here; \cref{maps:src:F.2.9}; \cref{maps:eq:F.31}}


No simply connected datum with $k=2$, $b_2=0$, and finite linear orders
$(3,5)$, $(3,8)$ or $(3,10)$ has the required maximal
square-zero cusp.  Among such data, the
remaining necessary unordered linear order pairs are
\[
 \{(3,4),(4,5),(5,6),(5,8),(5,12)\}.
 \tag{F.30}\label{maps:eq:F.30}
\]
When one order is five, the other finite characteristic
polynomial cannot be $(T-1)^2\Phi_4(T)^2$ or
$(T-1)^2\Phi_6(T)^2$.

\end{theorem}

\begin{proposition}[rational uniqueness of the extension]\label{maps:src:F.2.11}
\StatementRecord{maps:src:F.2.11}{proved}{Claude Fable 5.2 (Anthropic), effort high, 16 November 2026; also checked in this writing by \CurrentChecker}{the hypotheses, notation and definitions stated here; \cref{maps:src:F.2.9}; \cref{maps:eq:F.32}}


Let $W$ be a rational rank-five representation of a
twice-punctured affine line, with two finite generators
having fixed rank one and product with unipotent Jordan type
$(2,2,1)$.  Suppose $W$ has neither invariants nor
coinvariants.  Up to rescaling the invariant line, there is
exactly one rational extension
$0\to\mathbb Qe\to V\to W\to0$ in which the two finite
generators retain their finite orders and the cusp has
square-zero logarithm of rank three.  Thus, for the
two-point candidates of \ref{maps:src:F.2.9}, the remaining rational
representation problem is exactly the problem of the
rank-five quotients.  This statement does not identify
integral lattices in those representations.

\end{proposition}

\begin{theorem}[an integral obstruction for a rational $(5,12)$ family]\label{maps:src:F.2.12}
\StatementRecord{maps:src:F.2.12}{proved}{Claude Fable 5.2 (Anthropic), effort high, 16 November 2026; also checked in this writing by \CurrentChecker}{the hypotheses, notation and definitions stated here; \cref{maps:src:F.2.5}; \cref{maps:eq:F.35}}


Let a rational rank-six cohomology representation have generators
\[
 A=\begin{pmatrix}1&0\\0&A_5\end{pmatrix},\qquad
 B=\begin{pmatrix}1&c\\0&B_5\end{pmatrix},\qquad
 c=(0,0,0,1,0),
\]
where
\[
 A_5=\begin{pmatrix}
 1&0&1&0&0\\1&0&0&0&0\\-1&0&0&1&1\\
 1&-1&0&-1&0\\0&1&0&0&0
 \end{pmatrix},\qquad
 B_5=\begin{pmatrix}
 0&1&0&0&0\\0&0&0&0&1\\1&-2&0&0&0\\
 0&1&0&-1&-1\\0&0&1&1&2
 \end{pmatrix}.
 \tag{F.33}\label{maps:eq:F.33}
\]
For every invariant integral lattice with an index-one maximal
square-zero cusp, and every free affine filling of multiplicities
$m_1,m_2$ divisible respectively by five and twelve, a simply
connected total space has
\[
 \operatorname{Tor}_3 H_2
       \simeq\mathbb Z/3^{v_3(m_2)}.
 \tag{F.34}\label{maps:eq:F.34}
\]
In particular none of these data is three-connected.
The standard lattice nevertheless has a primitive cusp,
free affine fillings of orders five and twelve, and a simply
connected total space with rational Betti vector
$(1,0,0,0,4,0,0,0,1)$.

\end{theorem}

\begin{theorem}[two further rational families excluded integrally]\label{maps:src:F.2.13}
\StatementRecord{maps:src:F.2.13}{proved}{Claude Fable 5.2 (Anthropic), effort high, 16 November 2026; also checked in this writing by \CurrentChecker}{the hypotheses, notation and definitions stated here; \cref{maps:src:F.2.11}; \cref{maps:src:F.2.5}; \cref{maps:src:F.2.12}}


Let $C_f,C_g$ denote the degree-four companion matrices.
Take $f=\Phi_5$ and either
\[
 \text{case A: }g=\Phi_{12},\qquad
 \text{case B: }g=\Phi_3\Phi_6.
 \tag{F.36}\label{maps:eq:F.36}
\]
Let $W$ be the quotient of
$\bigwedge^2(C_f^{-1},C_g)$ by its invariant line, and
let $V$ be its unique rational extension of \ref{maps:src:F.2.11}.
For every invariant integral lattice in $V$ with primitive
maximal square-zero cusp, every simply connected free affine
filling, with multiplicities $m_1,m_2$ divisible by five and
six respectively, fails three-connectivity.
More precisely, put $r=v_3(m_2)\geq1$.
In case A, $\operatorname{Tor}_3 H_3$ has order
$3^{2r-1}$.  In case B,
$\operatorname{Tor}_3 H_2\simeq\mathbb Z/3^r$.
For $m_2=6$ the respective obstructing group is $\mathbb Z/3$.

\end{theorem}

\begin{theorem}[four order pairs and eight local spectral pairs]\label{maps:src:F.2.14}
\StatementRecord{maps:src:F.2.14}{proved}{Claude Fable 5.2 (Anthropic), effort high, 16 November 2026; also checked in this writing by \CurrentChecker}{the hypotheses, notation and definitions stated here; \cref{maps:src:F.2.15}; \cref{maps:src:F.2.9}; \cref{maps:src:F.2.7}; \cref{maps:eq:F.40}; \cref{maps:src:F.2.2}}


Every simply connected two-finite-point datum with $b_2=0$ has
one finite linear order equal to five and the other equal
to four, six, eight or twelve.
After suppressing the common factor $(T-1)^2$, its unordered
pair of finite characteristic polynomials is exactly one
of the following necessary possibilities:
\[
\begin{array}{c|c|l}
 \text{orders}&\text{order-five factor}&\text{other factor}\\\hline
 (5,4)&\Phi_5&\Phi_2^2\Phi_4\\
 (5,6)&\Phi_5&\Phi_2^2\Phi_3,\ \Phi_2^2\Phi_6,\ \Phi_3\Phi_6\\
 (5,8)&\Phi_5&\Phi_8\\
 (5,12)&\Phi_5&\Phi_{12},\ \Phi_3\Phi_4,\ \Phi_4\Phi_6.
\end{array}
 \tag{F.39}\label{maps:eq:F.39}
\]
Its rank-five quotient is rationally irreducible.
In particular all order-three pairs, including $(3,4)$,
are impossible.  The table is a necessary spectral filter;
it does not identify the relative rational or integral
positions of the matrices having those polynomials.

\end{theorem}

\begin{proposition}[automatic local complex compatibility]\label{maps:src:F.2.15}
\StatementRecord{maps:src:F.2.15}{proved}{Claude Fable 5.2 (Anthropic), effort high, 16 November 2026; also checked in this writing by \CurrentChecker}{the hypotheses, notation and definitions stated here; \cref{maps:src:F.2.7}; \cref{maps:src:F.1.4}}


Every simply connected datum of F.2 with two finite affine
fillings automatically has finite linear actions commuting
with real complex structures.  This assertion concerns the
individual finite actions and does not produce simultaneous
global periods.

\end{proposition}

\begin{theorem}[an unbounded integral family]\label{maps:src:F.2.16}
\StatementRecord{maps:src:F.2.16}{proved}{Claude Fable 5.2 (Anthropic), effort high, 16 November 2026; also checked in this writing by \CurrentChecker}{the hypotheses, notation and definitions stated here; \cref{maps:eq:F.U1}; \cref{maps:eq:F.U6}; \cref{maps:src:F.2.5}; \cref{maps:eq:F.U8}}


There are infinitely many simultaneous rational conjugacy classes
of integral rank-six tuples of the linear type $(5,4)$ in F.2,
with primitive maximal square-zero cusp.  An explicit infinite
subfamily has free affine lifts of orders five and four.
Apart from the parameter zero, the rank-five irreducible quotient
in the family is not self-dual.  Thus no finite bound on rational
conjugacy classes is valid for the unrestricted input tuples.
This does not assert infinitely many three-connected survivors.

\end{theorem}

\begin{proposition}[an odd torsion obstruction in this lattice]\label{maps:src:F.2.17}
\StatementRecord{maps:src:F.2.17}{proved}{Claude Fable 5.2 (Anthropic), effort high, 16 November 2026; also checked in this writing by \CurrentChecker}{the hypotheses, notation and definitions stated here; \cref{maps:eq:F.U6}; \cref{maps:src:F.2.19}; \cref{maps:eq:F.U5}; \cref{maps:src:F.2.12}}


For the integral tuple $(A_v,B_v)$ in \eqref{maps:eq:F.U6}, any simply
connected filling with free finite affine caps and the fixed
primitive Kuhn cusp has a subgroup
\[
 \mathbb Z/(v^2-v+1)\ \subset\operatorname{Tor}H_2(M;\mathbb Z).
 \tag{F.U9}\label{maps:eq:F.U9}
\]
In particular, in this specified family of lattices only
$v=0$ and $v=1$ can be three-connected, regardless of the
affine multiplicities and twists.  The stronger exclusion for every stable lattice in each
rational class is proved in Theorem \ref{maps:src:F.2.19} below.

\end{proposition}

\begin{theorem}[a survivor with a non-self-dual quotient]\label{maps:src:F.2.18}
\StatementRecord{maps:src:F.2.18}{proved}{Claude Fable 5.2 (Anthropic), effort high, 16 November 2026; also checked in this writing by \CurrentChecker}{the hypotheses, notation and definitions stated here; \cref{maps:eq:F.U6}; \cref{maps:src:F.3.11}; \cref{maps:src:F.2.5}; \cref{maps:src:F.2.16}; \cref{maps:src:F.2.6}; \cref{maps:eq:F.6}; \cref{maps:src:F.4.1}; \cref{maps:eq:F.U8}}


The parameter $v=1$ in \eqref{maps:eq:F.U6}, with affine orders five and four
and twists
\[
 w_1=(1,3,-2,-1,2,0)^t,\qquad
 w_2=(-1,0,0,0,0,0)^t,
\]
gives a three-connected smooth eight-manifold with
$H_4=\mathbb Z^4$ and all other reduced homology zero except
$H_8=\mathbb Z$.  It is almost complex and is oriented
diffeomorphic to $3\mathbb{HP}^2\#(-\mathbb{HP}^2)$.
Its irreducible rank-five quotient has no invariant bilinear
form, so integral three-connectivity does not force that
quotient to be orthogonal.  It has bounded-cusp holomorphic
periods by Theorem \ref{maps:src:F.3.11}; simultaneous finite holomorphic
extensions are not asserted here.

\end{theorem}

\begin{theorem}[a residual obstruction for every lattice]\label{maps:src:F.2.19}
\StatementRecord{maps:src:F.2.19}{proved}{Claude Fable 5.2 (Anthropic), effort high, 16 November 2026; also checked in this writing by \CurrentChecker}{the hypotheses, notation and definitions stated here; \cref{maps:src:F.2.16}; \cref{maps:eq:F.U6}; \cref{maps:src:F.2.6}; \cref{maps:src:F.2.18}; \cref{maps:eq:F.U12}; \cref{maps:src:F.2.17}; \cref{maps:eq:F.U1}}


Let a rational rank-six datum have finite linear orders five
and four, a primitive maximal cusp on a stable integral lattice,
and invariant line $L$, with irreducible rank-five quotient $W$
and no rational global coinvariants.  Suppose that at a prime
$p\ne2,5$ the semisimplified reduction of a stable lattice in
$W$ is the sum of the trivial module and an absolutely
irreducible module of dimension four.  Then no stable integral
lattice in that rational datum, with primitive maximal cusp,
and no free finite affine twists can give
$\pi_1=1$ and $H_2=0$ simultaneously.

For every $v\notin\{0,1\}$, the rational family in \ref{maps:src:F.2.16}
satisfies this hypothesis at each prime dividing
$v^2-v+1$.  Consequently the exclusion of those parameters
holds for \emph{all} stable lattices and all twists, not just
the lattice in \eqref{maps:eq:F.U6}.  The two remaining parameters are the
verified survivors of \ref{maps:src:F.2.6} and \ref{maps:src:F.2.18}.

\end{theorem}

\begin{theorem}[the complete rational $(5,4)$ normal form]\label{maps:src:F.2.20}
\StatementRecord{maps:src:F.2.20}{proved}{Claude Fable 5.2 (Anthropic), effort high, 16 November 2026; also checked in this writing by \CurrentChecker}{the hypotheses, notation and definitions stated here; \cref{maps:src:F.2.11}; \cref{maps:eq:F.R4}; \cref{maps:eq:F.R6}; \cref{maps:eq:F.R7}; \cref{maps:eq:F.R8}}


Consider ordered rational rank-five pairs $(A,C)$ with
\[
 \chi_A=(x-1)\Phi_5,
 \qquad \chi_C=(x-1)(x+1)^2(x^2+1).
\]
Put
\[
 N=CA^{-1}-I,
 \qquad N^2=0,\quad\operatorname{rank}N=2,
\]
where $A^5=C^4=I$ and the pair is irreducible over $\mathbb Q$.
Their simultaneous rational conjugacy classes are parametrized
bijectively by
\[
 \bigl\{[r:s:t]\in\mathbb P^2(\mathbb Q):D(r,s,t)\ne0\bigr\}
 \ \sqcup\ \mathbb A^1(\mathbb Q),
 \tag{F.R1}\label{maps:eq:F.R1}
\]
where
\[
\begin{split}
 D={}&r^3-4r^2s-r^2t+7rs^2+rst-3rt^2\\
     &{}-7s^3+11s^2t-11st^2+7t^3.
\end{split}
 \tag{F.R2}\label{maps:eq:F.R2}
\]
The second summand lies over the single exceptional point
$[4:1:3]$ of $D=0$.  The construction below is explicit and
uses only a ten-variable linear system.  By \ref{maps:src:F.2.11} it also
parametrizes the rational rank-six maximal-cusp extensions
of all these pairs, with $B=C^{-1}$.
This is a rational classification; the stable integral
lattices, affine twists and surviving homology still require
the arithmetic conditions of F.2.

\end{theorem}

\begin{corollary}[all good-prime lattices are forced]\label{maps:src:F.2.21}
\StatementRecord{maps:src:F.2.21}{proved}{Claude Fable 5.2 (Anthropic), effort high, 16 November 2026; also checked in this writing by \CurrentChecker}{the hypotheses, notation and definitions stated here; \cref{maps:src:F.2.20}; \cref{maps:src:F.2.19}}


For a $(5,4)$ datum with $\pi_1=1$ and $H_2=0$, the
rank-five quotient has absolutely irreducible reduction
at every prime $p\ne2,5$.  At each such prime its stable
$\mathbb Z_p$-lattice is unique up to homothety.
The same is true of the rank-six extension among lattices
with primitive maximal cusp.  Thus, once a rational pair
in \ref{maps:src:F.2.20} admits one survivor, all remaining integral-lattice
ambiguity is supported at the primes two and five.
This assertion supplies no finite bound on those bad-prime
lattices or on the rational parameters.

\end{corollary}

\begin{theorem}[a terminating lattice algorithm for each rational pair]\label{maps:src:F.2.22}
\StatementRecord{maps:src:F.2.22}{proved}{Claude Fable 5.2 (Anthropic), effort high, 16 November 2026; also checked in this writing by \CurrentChecker}{the hypotheses, notation and definitions stated here; \cref{maps:src:F.2.14}; \cref{maps:eq:F.L3}; \cref{maps:eq:F.L5}; \cref{maps:src:F.2.11}; \cref{maps:eq:F.L6}; \cref{maps:eq:F.L7}}


Let $(A,C)$ be any irreducible rational rank-five pair in
one of the eight spectral cases of \ref{maps:src:F.2.14}, with
$A^5=I$, $B=C^{-1}$, and $(AB-I)^2=0$ of rank two.
There is an explicit terminating algorithm which decides
whether the pair preserves an integral lattice, and, if so,
lists all such lattices up to simultaneous integral conjugacy.
It then lists all rank-six extensions with the same finite
orders and primitive maximal square-zero cusp, up to
simultaneous integral conjugacy.  The bounds below depend
on the input rational pair and are proved to lose no lattice.
They do not bound the rational parameter space or classify
all affine twists and resulting homology.

\end{theorem}

\begin{theorem}[complete lattice classes for the two family survivors]\label{maps:src:F.2.23}
\StatementRecord{maps:src:F.2.23}{proved}{Claude Fable 5.2 (Anthropic), effort high, 16 November 2026; also checked in this writing by \CurrentChecker}{the hypotheses, notation and definitions stated here; \cref{maps:src:F.2.16}; \cref{maps:src:F.2.6}; \cref{maps:src:F.2.18}; \cref{maps:eq:F.U1}; \cref{maps:eq:F.L4}; \cref{maps:eq:F.L5}; \cref{maps:eq:F.L6}; \cref{maps:eq:F.L7}; \cref{maps:src:F.2.12}; \cref{maps:eq:F.L10}; \cref{maps:eq:F.U6}; \cref{maps:src:F.2.19}}


For each of the two rational representations $v=0$ and $v=1$
of \ref{maps:src:F.2.16}, there are exactly two homothety classes of stable
rank-five quotient lattices, and exactly three integral
conjugacy classes of rank-six extensions with primitive
maximal cusp.  Exactly one of those three can have simply
connected finite affine fillings.  It is the integral class
already realized in \ref{maps:src:F.2.6}, respectively \ref{maps:src:F.2.18}.
Thus arbitrary stable-lattice changes produce no further
linear survivor classes anywhere in the family \ref{maps:src:F.2.16}.
This assertion retains the classification of all affine
twists and multiplicities as a separate problem.

\end{theorem}

\begin{theorem}[every surviving twist in the unbounded family]\label{maps:src:F.2.24}
\StatementRecord{maps:src:F.2.24}{proved}{Claude Fable 5.2 (Anthropic), effort high, 16 November 2026; also checked in this writing by \CurrentChecker}{the hypotheses, notation and definitions stated here; \cref{maps:src:F.2.12}; \cref{maps:src:F.2.16}; \cref{maps:src:F.2.23}; \cref{maps:src:F.2.20}; \cref{maps:src:F.2.19}; \cref{maps:eq:F.U6}; \cref{maps:src:F.2.5}; \cref{maps:src:F.2.18}; \cref{maps:eq:F.T4}; \cref{maps:src:F.2.6}; \cref{maps:eq:F.T6}; \cref{maps:eq:F.T7}; \cref{maps:eq:F.T5}; \cref{maps:eq:F.T8}; \cref{maps:eq:F.T3}; \cref{maps:eq:F.T2}; \cref{maps:eq:F.T9}}


Use the fixed Kuhn cusp and the affine translation clutching
convention of \ref{maps:src:F.2.12}.  In the entire rational family \ref{maps:src:F.2.16},
allow every stable lattice, every free affine finite lift,
and every pair of actual finite multiplicities.
The data with $\pi_1=1$ and $H_2=H_3=0$ are precisely the
following, up to the integral linear classification already
proved and the translation identifications described below:
\[
\begin{gathered}
 v\in\{0,1\},\qquad
 \text{the unique eligible integral class of \ref{maps:src:F.2.23}},\\
 (m_1,m_2)=(5,4),\qquad
 20\kappa+4r_1+5r_2=\pm1.
\end{gathered}
 \tag{F.T1}\label{maps:eq:F.T1}
\]
Here $r_i=\ell(v_i)$ is the invariant affine numerator, and
$\kappa\in\mathbb Z$ is the invariant coordinate of the
cusp translation clutching, with signs fixed by the displayed
van Kampen relation below.  The remaining fixed-direction
components of the twists impose no condition.
Every manifold in \eqref{maps:eq:F.T1} is oriented diffeomorphic to its
corresponding known survivor and hence to
$3\mathbb{HP}^2\#(-\mathbb{HP}^2)$.
This completes the survivor classification for this entire
unbounded rational family; it does not classify the other
rational points of \ref{maps:src:F.2.20} or the other spectral loci.

\end{theorem}

\begin{theorem}[the universal degree-two suspension obstruction]\label{maps:src:F.2.25}
\StatementRecord{maps:src:F.2.25}{proved}{Claude Fable 5.2 (Anthropic), effort high, 16 November 2026; also checked in this writing by \CurrentChecker}{the hypotheses, notation and definitions stated here; \cref{maps:src:F.2.14}; \cref{maps:src:F.2.7}; \cref{maps:src:F.2.9}; \cref{maps:src:F.2.17}; \cref{maps:src:F.2.12}; \cref{maps:eq:F.T4}; \cref{maps:eq:F.S3}; \cref{maps:eq:F.S6}; \ExternalPaperRef{theorem~}{6.3}{III}; \cref{maps:eq:F.S2}; \cref{maps:eq:F.S7}; \cref{maps:eq:F.S5}; \cref{maps:eq:F.S8}}


Every datum of F.2 with $\pi_1=1$ and $H_2=H_3=0$ has
exactly two genuine finite fillings and has minimal affine
multiplicities.  After ordering the order-five point first,
its finite characteristic polynomials are one of precisely
the following two remaining necessary possibilities:
\[
\begin{array}{c|c|c}
 (m_1,m_2)&\chi_{M_1}&\chi_{M_2}\\\hline
 (5,4)&(x-1)^2\Phi_5&(x-1)^2\Phi_2^2\Phi_4\\
 (5,6)&(x-1)^2\Phi_5&(x-1)^2\Phi_2^2\Phi_3.
\end{array}
 \tag{F.S1}\label{maps:eq:F.S1}
\]
In particular all order-eight and order-twelve spectral
loci, and the other two order-six spectral loci of \ref{maps:src:F.2.14},
are excluded for every rational relative position, every
stable lattice and every free affine twist.
There is no multiplicity search bound in this statement:
the only possible actual affine multiplicities are the
two displayed pairs.

\end{theorem}

\begin{theorem}[the complete rational remaining $(5,6)$ locus]\label{maps:src:F.2.26}
\StatementRecord{maps:src:F.2.26}{proved}{Claude Fable 5.2 (Anthropic), effort high, 16 November 2026; also checked in this writing by \CurrentChecker}{the hypotheses, notation and definitions stated here; \cref{maps:eq:F.S1}; \cref{maps:eq:F.R2}; \cref{maps:src:F.2.25}; \cref{maps:eq:F.R3}; \cref{maps:eq:F.R5}; \cref{maps:src:F.2.20}; \cref{maps:eq:F.R6}; \cref{maps:eq:F.R7}; \cref{maps:eq:F.R8}; \cref{maps:eq:F.R12}; \cref{maps:src:F.2.22}}


The irreducible rational rank-five pairs in the second row
of \eqref{maps:eq:F.S1}, with square-zero rank-two cusp, are classified
up to simultaneous rational conjugacy by
\[
 \{[r:s:t]\in\mathbb P^2(\mathbb Q):
             D(r,s,t)\ne0,\ Q_-(r,s,t)\ne0\}
       \ \sqcup\ \mathbb A^1(\mathbb Q),
 \qquad Q_-=r^2-rs+st,
 \tag{F.R10}\label{maps:eq:F.R10}
\]
where $D$ is the same cubic as in \eqref{maps:eq:F.R2}.
The exceptional line again lies over $[4:1:3]$.
Thus the two rational spectral loci left by \ref{maps:src:F.2.25} both
have complete explicit normal forms, including their
entire rational exceptional loci.

\end{theorem}

\begin{theorem}[the entire remaining order-six locus is excluded]\label{maps:src:F.2.27}
\StatementRecord{maps:src:F.2.27}{proved}{Claude Fable 5.2 (Anthropic), effort high, 16 November 2026; also checked in this writing by \CurrentChecker}{the hypotheses, notation and definitions stated here; \cref{maps:eq:F.S1}; \cref{maps:src:F.2.25}; \cref{maps:eq:F.S3}; \cref{maps:eq:F.S5}; \cref{maps:eq:F.X2}; \cref{maps:eq:F.X3}; \cref{maps:eq:F.S7}; \cref{maps:eq:F.X5}; \ExternalPaperRef{equation~}{{E.9}}{III}}


Every simply connected datum in the second row of \eqref{maps:eq:F.S1},
at every rational relative position and every stable integral
lattice with primitive maximal cusp, has
\[
                 \operatorname{Tor}_3 H_3(M;\mathbb Z)\ne0
\]
if its degree-two suspension cokernels vanish.
Consequently no such datum is three-connected.  Together with
\ref{maps:src:F.2.25} this leaves only the minimal order pair $(5,4)$ as a
possible three-connected datum in the full problem F.2.
\end{theorem}

\begin{theorem}[universal fundamental group and twist normalization]\label{maps:src:F.2.28}
\StatementRecord{maps:src:F.2.28}{proved}{Claude Fable 5.2 (Anthropic), effort high, 16 November 2026; also checked in this writing by \CurrentChecker}{the hypotheses, notation and definitions stated here; \cref{maps:src:F.2.25}; \cref{maps:eq:F.S5}; \cref{maps:eq:F.T8}; \cref{maps:src:F.2.24}}


In the remaining $(5,4)$ locus, fix the marked standard Kuhn
cap and allow affine translation clutching.
Let $\Lambda=\mathbb Z^6$ be fibre homology and let
$\Lambda_\Gamma$ be its global monodromy coinvariants.
Write $\Lambda_\Gamma=\mathbb Zh\oplus T$, with $T$ finite,
and let $\ell$ be the primitive invariant covector detecting
the free factor.  The following are complete reductions for
any datum that could satisfy $\pi_1=1$ and $H_2=H_3=0$.

The actual multiplicities must be $(5,4)$, and the local
vertical fixed pairings must be $h_1=5,h_2=4$, as in \ref{maps:src:F.2.25}.
For any affine twists and clutching the fundamental group
is abelian.  If its free presentation determinant is nonzero,
its order is
\[
 |\pi_1|=|20\kappa+4r_1+5r_2|\,|T|,
 \qquad r_i=\ell(v_i).
 \tag{F.T11}\label{maps:eq:F.T11}
\]
If that determinant is zero the group is infinite.
In particular simple connectivity is equivalent to
\[
 T=0,\qquad 20\kappa+4r_1+5r_2=\pm1.
 \tag{F.T12}\label{maps:eq:F.T12}
\]
Whenever these conditions and $h_1=5,h_2=4$ hold, every
such twist choice is oriented diffeomorphic to the normalized
filling with twists $(u_1,-u_2)$ and zero cusp clutching,
where $u_i$ is any integral fixed vector with $\ell(u_i)=1$.
Thus each eligible integral linear datum has exactly one
normalized smooth filling to test for the remaining
homology conditions; no affine-twist enumeration remains.

\end{theorem}

\begin{theorem}[a complete finite integral survivor test]\label{maps:src:F.2.29}
\StatementRecord{maps:src:F.2.29}{proved}{Claude Fable 5.2 (Anthropic), effort high, 16 November 2026; also checked in this writing by \CurrentChecker}{the hypotheses, notation and definitions stated here; \cref{maps:src:F.2.20}; \cref{maps:src:F.2.22}; \cref{maps:src:F.2.25}; \cref{maps:src:F.2.28}; \cref{maps:src:F.2.6}; \cref{maps:eq:F.HC7}; \ExternalPaperRef{theorem~}{6.3}{III}; \cref{maps:src:F.2.16}}


Take any rational pair in \ref{maps:src:F.2.20} and any of its integral,
primitive-cusp rank-six extensions supplied by \ref{maps:src:F.2.22}.
Write $M_1,M_2$ for geometric transport on fibre homology
$\Lambda=\mathbb Z^6$, in the peripheral convention
$M_0M_1M_2=I$.
First require the preliminary conditions
\[
 \Lambda_\Gamma=\mathbb Z,\qquad
 \ell(\Lambda^{M_i})=\mathbb Z,\qquad h_1=5,\quad h_2=4 .
 \tag{F.HC1}\label{maps:eq:F.HC1}
\]
Here $\ell$ is the primitive global invariant covector, and
$h_i$ is the integral fixed-vector/fixed-covector pairing
on $\ker\ell$ from \ref{maps:src:F.2.25}.
Failure of any condition excludes the datum.
If they hold, use the unique normalized affine filling of
\ref{maps:src:F.2.28}, with twists $(u_1,-u_2)$ and zero cusp clutching.
The following six finite integer-matrix tests are necessary
and sufficient for this filling to have $H_2=H_3=0$.

Let
\[
 T_i=M_i^{-t},\quad T_0=(T_1T_2)^{-1},\qquad
 G_q=\operatorname{coker}
 [\,\bigwedge^qT_1-I\mid\bigwedge^qT_2-I\,]
 \quad(1\le q\le4).
 \tag{F.HC2}\label{maps:eq:F.HC2}
\]
All $G_q$ are finite.  Let $P_q$ be the integral parabolic
cohomology group of the exterior local system
$\bigwedge^q\mathcal V$, for $q=2,4$.
Let $t_q\in G_{q-1}$, $2\le q\le5$, be the natural global
coinvariant image of the unique nonzero torsion suspension
class at the order-four cap, as specified below.
Let $c_4\in G_3$ be the natural image of the unique nonzero
torsion class in
$\operatorname{coker}(\bigwedge^3T_0-I)$.
Then the six tests are
\[
\begin{array}{c|ccc}
 &\text{first test}&\text{second test}&\text{third test}\\ \hline
 H^3=0&G_1=\langle t_2\rangle&P_2=0&t_3\ne0\\
 H^5=0&G_3=\langle t_4,c_4\rangle&P_4=0&t_5\ne0 .
\end{array}
\tag{F.HC3}\label{maps:eq:F.HC3}
\]
Every group, class, and equality here is computed by
finitely many integer kernels, Smith forms, and exterior
minors.  This is a decision criterion for every rational
parameter and every integral lattice in the complete
normal form; it makes no bounded-height assertion about
the rational parameter set.

\end{theorem}

\begin{corollary}[constructive classification of all survivors]\label{maps:src:F.2.30}
\StatementRecord{maps:src:F.2.30}{proved}{Claude Fable 5.2 (Anthropic), effort high, 16 November 2026; also checked in this writing by \CurrentChecker}{the hypotheses, notation and definitions stated here; \cref{maps:src:F.2.20}; \cref{maps:src:F.2.22}; \cref{maps:eq:F.HC1}; \cref{maps:eq:F.HC3}; \cref{maps:eq:F.T12}; \cref{maps:src:F.3.11}; \cref{maps:src:F.2.7}; \cref{maps:src:F.2.14}; \cref{maps:src:F.2.25}; \cref{maps:src:F.2.27}; \cref{maps:src:F.2.28}; \cref{maps:src:F.2.29}; \ExternalPaperRef{theorem~}{6.3}{III}; \cref{maps:src:F.2.15}; \cref{maps:src:F.2.5}; \S\ref{sec:V-known-inputs}; \cref{maps:src:F.4.1}}


For the marked Kuhn cap and affine translation clutchings
used throughout F.2, the following is a complete constructive
classification of the data with
$\pi_1=1$ and $H_2=H_3=0$, up to integral conjugacy of their
linear tuples and oriented diffeomorphism of their fillings.

There are no data with three genuine finite fillings.
With two fillings their actual affine orders are five and
four.  Order the order-five point first.
The complete, unbounded parameter set is the rational
projective normal form of \ref{maps:src:F.2.20} together with its
exceptional rational affine line.
For each parameter, run the terminating order,
conductor, stable-lattice, and primitive-extension algorithm
of \ref{maps:src:F.2.22}.  Discard precisely the integral classes failing
\eqref{maps:eq:F.HC1} or one of \eqref{maps:eq:F.HC3}.
The classes retained by this procedure are exactly all
surviving integral linear tuples.
For each retained class all simply connected affine
twists are exactly the choices satisfying \eqref{maps:eq:F.T12};
their fillings are all oriented diffeomorphic to its
normalized filling.
If ordered labels with the order-four point first are
desired, interchange the finite labels and recompute
$M_0=(M_1M_2)^{-1}$.

Every retained datum has
\[
 \pi_1=1,\qquad
 H_j=\begin{cases}
  \mathbb Z,&j=0,8,\\
  \mathbb Z^4,&j=4,\\
  0,&\text{otherwise},
 \end{cases}
 \tag{F.HC13}\label{maps:eq:F.HC13}
\]
and its oriented smooth manifold is
$3\mathbb{HP}^2\#(-\mathbb{HP}^2)$.
It has bounded-cusp holomorphic periods in the literal
sense of F.3, by \ref{maps:src:F.3.11}.
The classification here concerns the survivor subproblem.
It does not classify the nonsurviving tuples or compute
their full integral homology, as also requested in F.2.

\end{corollary}

\begin{theorem}[complete fundamental-group data for every input]\label{maps:src:F.2.31}
\StatementRecord{maps:src:F.2.31}{proved}{Claude Fable 5.2 (Anthropic), effort high, 16 November 2026; also checked in this writing by \CurrentChecker}{the hypotheses, notation and definitions stated here; \ExternalPaperRef{theorem~}{6.3}{III}; \cite[Theorems 2.3--2.4, p.~425; Euler-characteristic comparison, pp.~427--428]{rootScott83}}


Let $k=2$ or $3$, and let an integral rank-six tuple satisfy the
finite-order and primitive maximal-cusp conditions of F.2.  No
irreducibility, determinant-one, or prospective-survivor hypothesis is
imposed.  Specify each actual finite affine action as
\[
 \tau_i(x)=M_ix+v_i/m_i,\qquad v_i\in\Lambda=\mathbb Z^6,
 \quad \tau_i^{m_i}=1\text{ on }\Lambda_{\mathbb R}/\Lambda,
 \tag{F.G1}\label{maps:eq:F.G1}
\]
and require its cyclic action to be free.  The numerator need not be
fixed by $M_i$.  Define its actual integral power vector and the cusp
lattice by
\[
 w_i=\frac1{m_i}\sum_{j=0}^{m_i-1}M_i^jv_i\in\Lambda^{M_i},
 \qquad U=\operatorname{im}\bigl((M_1\cdots M_k)^{-1}-I\bigr).
 \tag{F.G2}\label{maps:eq:F.G2}
\]
Use left conjugation transport for the peripheral generators $s_i$,
and record the integral translation clutching by $\kappa\in\Lambda$
in $s_1\cdots s_kx^\kappa=1$.  Then the full fundamental group is
\[
 \left\langle x^\lambda,s_1,\ldots,s_k\ \left|\
 \begin{array}{l}
 x^\lambda x^\mu=x^{\lambda+\mu},\quad x^u=1\ (u\in U),\\
 s_ix^\lambda s_i^{-1}=x^{M_i\lambda},\quad
 s_i^{m_i}=x^{w_i},\quad s_1\cdots s_kx^\kappa=1
 \end{array}\right.\right\rangle .                         \tag{F.G3}\label{maps:eq:F.G3}
\]
Only a basis of $\Lambda$ and of $U$ is needed in this finite
presentation.  The program below also determines exactly which
quotient of the apparent fibre subgroup injects.

Put
\[
 L=\sum_{g\in\langle M_1,\ldots,M_k\rangle}gU,
 \qquad A=\Lambda/L.                                      \tag{F.G4}\label{maps:eq:F.G4}
\]
This is an integral sum, without saturation.  For two points, put
$d=\gcd(m_1,m_2)$, $M=M_1$, and
$S_j(M)=I+M+\cdots+M^{j-1}$.  Choose integers $a,b$ with
$am_1+bm_2=d$, and set, in $A$,
\[
 \begin{split}
 w'_2&=-w_2-S_{m_2}(M)\kappa,\\
 e&=(m_2/d)w_1-(m_1/d)w'_2,\qquad
 z=aw_1+bw'_2,\qquad \overline A=A/\langle e\rangle .
 \end{split}                                             \tag{F.G5}\label{maps:eq:F.G5}
\]
There is an exact sequence and a complete marked extension description
\[
 1\longrightarrow\overline A\longrightarrow\pi_1
   \longrightarrow C_d\longrightarrow1,
 \qquad txt^{-1}=M(x),\quad t^d=z.                         \tag{F.G6}\label{maps:eq:F.G6}
\]
Thus its order is $d|\overline A|$ if $\overline A$ is finite;
otherwise it is infinite.  In particular $d=1$ gives precisely
the abelian group $\overline A$.

For three genuine points, so $m_i\ge2$, let
\[
 \Gamma=\langle q_1,q_2,q_3\mid q_i^{m_i}=1,
                         \ q_1q_2q_3=1\rangle,
 \qquad \epsilon=\sum_i1/m_i-1.                           \tag{F.G7}\label{maps:eq:F.G7}
\]
The action $q_i\mapsto M_i$ is well defined on $A$.
If $\epsilon\le0$, the exact sequence is
\[
                    1\longrightarrow A\longrightarrow\pi_1
                         \longrightarrow\Gamma\longrightarrow1. 
 \tag{F.G8}\label{maps:eq:F.G8}
\]
If $\epsilon>0$, $\Gamma$ is the finite spherical triangle group.
Using its action on $A$, form the exact transfer vector
\[
 e=\sum_{g\in\Gamma}g\kappa+
       \sum_{i=1}^3\ \sum_{g\in\Gamma/\langle q_i\rangle}g w_i.
 \tag{F.G9}\label{maps:eq:F.G9}
\]
The second sums use left coset representatives and are independent
of their choice.  In this case the exact sequence has kernel
$A/\langle e\rangle$, with the same lifted relations \eqref{maps:eq:F.G3}.
Consequently \eqref{maps:eq:F.G3}, \eqref{maps:eq:F.G6}, \eqref{maps:eq:F.G8}, and \eqref{maps:eq:F.G9} give every requested
fundamental group by finite operations, including the precise
injected abelian kernel and all extension data.  These statements
do not claim a conjugacy classification of the original tuples.

For every input, the full first integral homology is the Smith
cokernel of the integer matrix
\[
 \begin{pmatrix}
 M_1-I&\cdots&M_k-I&-w_1&\cdots&-w_k&\kappa\\
 0&\cdots&0&&\operatorname{diag}(m_1,\ldots,m_k)&&\boldsymbol1
 \end{pmatrix}.                                           \tag{F.G10}\label{maps:eq:F.G10}
\]
Here the bottom left zero block has $k$ rows.  The code records
the full matrix, not only its invariant factors.

\end{theorem}

\begin{theorem}[a complete three-point logarithmic family]\label{maps:src:F.2.32}
\StatementRecord{maps:src:F.2.32}{proved}{Claude Fable 5.2 (Anthropic), effort high, 16 November 2026; also checked in this writing by \CurrentChecker}{the hypotheses, notation and definitions stated here; \cref{maps:src:F.3.7}; \cref{maps:src:F.2.28}; \cref{maps:eq:F.LT4}; \cref{maps:eq:F.LT5}; \cref{maps:src:F.2.5}; \cref{maps:eq:F.HC8}; \cref{maps:eq:F.HC9}; \cref{maps:eq:F.LT8}; \cref{maps:src:F.2.29}; \cref{maps:eq:F.LT6}; \cref{maps:eq:F.LT9}}


Start with the compact complex fourfold of \ref{maps:src:F.3.7} and its
normalized order-five and order-four affine fillings.
Let $\ell$ be the primitive invariant covector on fibre
homology, with the sign chosen in \ref{maps:src:F.2.28}.
At any additional ordinary fibre choose
$v\in\mathbb Z^6$ and an integer $m\ge2$ such that
\[
 r=\ell(v),\qquad \gcd(r,m)=1.
\tag{F.LT1}\label{maps:eq:F.LT1}
\]
Perform the marked logarithmic transformation with
translation $v/m$ described below.
This gives a smooth compact complex fourfold $X_{m,v}$
with Euler number six and three genuine finite affine
fillings, whose linear orders are five, four and one
and whose actual affine orders are five, four and $m$.
Its fundamental group is
\[
 \pi_1(X_{m,v})=
 \langle a,b,c,h\mid h\text{ central},\
   a^5=h,\ b^4=h^{-1},\ c^m=h^r,\ abc=1\rangle .
\tag{F.LT2}\label{maps:eq:F.LT2}
\]
This is the group of the oriented Seifert manifold
with invariants
$(0;(5,-1),(4,1),(m,-r))$, in the convention
$x_i^{\alpha_i}h^{\beta_i}=1$ and product $x_i=1$.
Its base orbifold is hyperbolic, and its central fibre
$h$ has infinite order.

Put $C_{m,r}=\mathbb Z/(20r-m)\mathbb Z$, with
$\mathbb Z/0\mathbb Z=\mathbb Z$, and let
$\varepsilon=1$ when $20r=m$ and zero otherwise.
All the integral homology groups are
\[
\begin{array}{c|c}
j&H_j(X_{m,v};\mathbb Z)\\ \hline
0,8&\mathbb Z\\
1,6&C_{m,r}\\
2&\mathbb Z^\varepsilon\oplus(\mathbb Z/m)^5\\
3&(\mathbb Z/m)^{10}\\
4&\mathbb Z^4\oplus(\mathbb Z/m)^{10}\\
5&(\mathbb Z/m)^5\\
7&\mathbb Z^\varepsilon .
\end{array}
\tag{F.LT3}\label{maps:eq:F.LT3}
\]
Under \eqref{maps:eq:F.LT1} the exceptional free-rank case is exactly
$(m,r)=(20,1)$.
In particular $(m,r)=(19,1)$ and $(21,1)$ have
$H_1=0$ but nontrivial, infinite fundamental groups.
More generally this occurs whenever $|20r-m|=1$.
The result treats every vertical component of $v$
in $\ker\ell$, not only a chosen fixed representative.

\end{theorem}

\begin{theorem}[arbitrary finite affine caps]\label{maps:src:F.2.33}
\StatementRecord{maps:src:F.2.33}{proved}{GPT-6 Astra (OpenAI, ChatGPT Pro mode), effort Pro, 2026-11-17 UTC; GPT-6 Astra (OpenAI, ChatGPT Pro mode, effort Pro),17November2026,/root; Claude Fable 5.2 (Anthropic), effort high, 16 November 2026; also checked in this writing by \CurrentChecker}{\cref{maps:src:F.2.31}; the displayed affine actions and bar-chain maps}


Let $g(x)=Mx+v/m$ be a free cyclic affine action of actual order
$m\geq2$ on $T^r$, with $M\in\mathrm{GL}(r,\mathbb Z)$ and
$v\in\mathbb Z^r$.  For every specified degree bound $d$, there is
an explicit terminating integer algorithm producing chains and
boundary maps through degree $d+1$ for the pair
\[
 B=(D^2\times T^r)/\langle(s,x)\mapsto
 (e^{2\pi i/m}s,g(x))\rangle,
 \qquad Y=\partial B.                                      \tag{F.CA1}\label{maps:eq:F.CA1}
\]
The cap complex has at most $2^r$ generators in each degree and
the boundary complex at most $2^{r+1}$.  These are complexes of
finite rank in each degree; only their indicated finite prefixes
are required to compute $H_j(B)$, $H_j(Y)$ and $H_j(B,Y)$ for
$j\leq d$.  Their maps retain the affine translations and all
integral higher coherence terms.  In particular, the construction
does not replace an equivariant torus chain complex by its
homology without recording the correction terms.

\end{theorem}

\begin{theorem}[the marked Kuhn cusp pair]\label{maps:src:F.2.34}
\StatementRecord{maps:src:F.2.34}{proved}{GPT-6 Astra (OpenAI, ChatGPT Pro mode), effort Pro, 2026-11-17 UTC; GPT-6 Astra (OpenAI, ChatGPT Pro mode, effort Pro),17November2026,/root; Claude Fable 5.2 (Anthropic), effort high, 16 November 2026; also checked in this writing by \CurrentChecker}{\cref{maps:src:F.2.33}; \cite[Theorem 4.1.1 and Lemmas 4.1.2--4.1.3, pp.~480--482]{ffAchingerOgus20}; \cite[Theorem 8.7, pp.~57--59, and Corollary 8.8, p.~59, attached source build]{manYang26}}


The fixed primitive rank-three Kuhn cusp admits an explicit
finite integer chain complex for its cap $W$, its boundary
$Y$, and the inclusion $Y\to W$.  It is obtained from a
category with 26 objects, 170 nonidentity morphisms, and
respectively $26,170,288,144$ strict flags of lengths
$0,1,2,3$.  The construction records the full integral
collapse map.  In particular it gives
\[
\begin{array}{c|rrrrrrrrr}
j&0&1&2&3&4&5&6&7&8\\\hline
\operatorname{rank}H_j(W)&1&3&9&9&10&3&1&0&0\\
\operatorname{rank}H_j(Y)&1&4&12&18&18&12&4&1&0\\
\operatorname{rank}H_j(W,Y)&0&0&1&3&10&9&9&3&1
\end{array}                                                \tag{F.CK1}\label{maps:eq:F.CK1}
\]
The only torsion in these three rows is
$\operatorname{Tor}H_3(Y)=\mathbb Z/2$.

\end{theorem}

\begin{theorem}[complete integral homology of a marked gluing]\label{maps:src:F.2.35}
\StatementRecord{maps:src:F.2.35}{proved}{GPT-6 Astra (OpenAI, ChatGPT Pro mode), effort Pro, 2026-11-17 UTC; GPT-6 Astra (OpenAI, ChatGPT Pro mode, effort Pro),17November2026,/root; Claude Fable 5.2 (Anthropic), effort high, 16 November 2026; also checked in this writing by \CurrentChecker}{\cref{maps:src:F.2.31}, \cref{maps:src:F.2.33}, \cref{maps:src:F.2.34}; the displayed complement and peripheral maps}


For every fully marked F.2 datum there is a terminating exact
integer algorithm computing every integral homology group of
the resulting eight-manifold, including every torsion invariant
factor.  The input consists of the finite homology matrices
$M_1,\ldots,M_k$, free affine numerators $v_i/m_i$ of their
actual orders, the integral clutch vector $\kappa$, and the
chosen primitive cusp marking
\[
 M_0=(M_1\cdots M_k)^{-1},\qquad
 M_0P=P\begin{pmatrix}I_3&I_3\\0&I_3\end{pmatrix},
 \qquad P\in\mathrm{GL}(6,\mathbb Z).                       \tag{F.GC1}\label{maps:eq:F.GC1}
\]
Here $k=2$ or $3$ and the genuine orders are at least two;
order-one sites are first eliminated as in \ref{maps:src:F.2.31}.
If a marking is not supplied, the algorithm makes and records
one by Smith reduction.  It does not identify that choice
with an unspecified prescribed marking.
The complete chain prefix through degree nine has at most
4003 generators in any degree.  No spectral-sequence
extension, saturation convention, orientation assumption or
unbounded search enters the homology calculation.

\end{theorem}

\begin{theorem}[simultaneous integral conjugacy for every input]\label{maps:src:F.2.36}
\StatementRecord{maps:src:F.2.36}{proved}{Claude Fable 5.2 (Anthropic), effort high, 16 November 2026; also checked in this writing by \CurrentChecker}{the hypotheses, notation and definitions stated here; \cite[Algorithm 8.3, p.~21]{rootBHJ22}; \cite[Algorithm 8.3 and its correctness proof, p.~21]{rootBHJ22}; \cite[Theorem 3.4, p.~6; Propositions 8.1--8.2, p.~21]{rootBHJ22}}


There is an unconditional exact algorithm which decides simultaneous
$\mathrm{GL}_6(\mathbb Z)$ conjugacy of any two marked matrix tuples of
F.2 and returns a conjugating matrix when one exists. The algorithm
does not impose a bound on their entries, on the affine multiplicities,
or on the integral conjugator. More generally it applies to any integral
tuple containing a matrix $1+N$ with $N^2=0$ and $\operatorname{rank}N=3$;
the primitive-image assumption is unnecessary for this statement.

\end{theorem}

\begin{theorem}[a complete canonical-representative algorithm]\label{maps:src:F.2.37}
\StatementRecord{maps:src:F.2.37}{proved}{GPT-6 Astra (OpenAI, ChatGPT Pro mode), effort Pro, 2026-11-17 UTC; GPT-6 Astra (OpenAI, ChatGPT Pro mode, effort Pro),17November2026,/root; Claude Fable 5.2 (Anthropic), effort high, 16 November 2026; also checked in this writing by \CurrentChecker}{the hypotheses, notation and definitions stated here; \cref{maps:src:F.2.38}; \cref{maps:eq:F.CN1}; \cref{maps:src:F.2.36}; \cref{maps:eq:F.CN2}; \cref{maps:src:F.2.16}}


Every marked linear tuple in F.2, for either $k=2$ or $k=3$,
has a computable canonical integral representative and a
returned simultaneous unimodular conjugator. Two inputs have
the same representative if and only if they are simultaneously
$\mathrm{GL}_6(\mathbb Z)$ conjugate. The algorithm uses a proved
finite bound depending on the input. No fixed global entry bound
can represent every class.

This is an effective classification, not a claim that its
infinite output has been printed in a finite table. The
ordered involution sector has the separate complete finite
normal form \ref{maps:src:F.2.38}, and all of its $56$ classes have been
computed. The general bounded fallback is potentially very
expensive; its finiteness is a theorem, not a runtime estimate.

\end{theorem}

\begin{theorem}[the complete ordered involution stratum]\label{maps:src:F.2.38}
\StatementRecord{maps:src:F.2.38}{proved}{GPT-6 Astra (OpenAI, ChatGPT Pro mode), effort Pro, 2026-11-17 UTC; GPT-6 Astra (OpenAI, ChatGPT Pro mode, effort Pro),17November2026,/root; Claude Fable 5.2 (Anthropic), effort high, 16 November 2026; also checked in this writing by \CurrentChecker}{the hypotheses, notation and definitions stated here; \cref{maps:eq:F.IV1}; \cref{maps:eq:F.IV3}; \cref{maps:eq:F.IV4}; \cref{maps:eq:F.IV2}; \cref{maps:src:F.2.31}}


There are exactly 56 simultaneous $\mathrm{GL}_6(\mathbb Z)$
conjugacy classes of ordered tuples $(M_1,M_2,M_0)$ satisfying
the primitive maximal-cusp conditions of F.2 with both finite
linear orders equal to two.  An explicit integral normalizer
returns a canonical representative and a checked unimodular
conjugating matrix for every such input, without any bound on
its entries.  Every representative has free affine lifts of
actual order four at both finite points.  None can give a
simply connected filling, for any choices of free affine
orders or twists.

\end{theorem}

\begin{theorem}[the full effective classification and all marked gluings]\label{maps:src:F.2.39}
\StatementRecord{maps:src:F.2.39}{proved}{GPT-6 Astra (OpenAI, ChatGPT Pro mode), effort Pro, 2026-11-17 UTC; GPT-6 Astra (OpenAI, ChatGPT Pro mode, effort Pro),17November2026,/root; Claude Fable 5.2 (Anthropic), effort high, 16 November 2026; also checked in this writing by \CurrentChecker}{the hypotheses, notation and definitions stated here; \cref{maps:src:F.2.37}; \cref{maps:src:F.2.38}; \cref{maps:src:F.2.31}; \cref{maps:src:F.2.35}; \cref{maps:src:F.2.3}; \cref{maps:src:F.2.30}; \cref{maps:src:F.2.1}; \cref{maps:src:F.2.2}; \cref{maps:src:F.2.36}; \cref{maps:eq:F.CN3}; \cref{maps:src:F.2.33}; \cref{maps:src:F.2.7}}


The linear classification and all-twist topological computation
asked for in F.2 have the following complete effective answer.
For $k=2$ and $k=3$, every admissible ordered integral tuple
has the canonical representative of \ref{maps:src:F.2.37}, with an explicit
simultaneous unimodular conjugator.  Equality of these
representatives is necessary and sufficient for simultaneous
integral conjugacy.  There is no fixed global height cutoff;
the finite stopping bound for an input of height $H$ is
$(2H+1)^{36(k+1)}$ candidates before the proved rejection
filters.  The separate involution normal form has exactly
56 classes and is computed completely in \ref{maps:src:F.2.38}.

For every choice of actual free affine orders and integer
numerators, integral clutch vector, and specified primitive
Kuhn cusp marking, \ref{maps:src:F.2.31} and \ref{maps:src:F.2.35} compute the full
fundamental group and every integral homology group of the
gluing.  This is the fully marked correction to \ref{maps:src:F.2.3};
peripheral matrices are not substituted for a choice of
filling.  The algorithms have no cutoff on affine orders,
twist entries or lattice indices.  The complete list of data
with $\pi_1=1$ and $H_2=H_3=0$ is the constructive
parametrization of \ref{maps:src:F.2.30}, with its exact lattice enumeration
and necessary-and-sufficient finite tests.  Each retained
filling has the homology and oriented smooth type stated
there.

\end{theorem}

\begin{theorem}\label{night:P1B.1}
\StatementRecord{night:P1B.1}{proved}{S6\_06\_BLIND: GPT-6 Astra (OpenAI, ChatGPT Pro mode, effort Pro), recorded re-derivation in the checkpoint dated 18 November 2026; also checked in this writing by \CurrentChecker}{Definitions and the displayed constructions}
Let $N$ be an endomorphism of $\Z^6$ with $N^2=0$, rank three,
and primitive image. There is a unimodular basis in which
\[
 I+N=\begin{pmatrix}I_3&I_3\\0&I_3\end{pmatrix}.
\]
Its centralizer in $\operatorname{GL}_6(\Z)$ consists exactly of
$\left(\begin{smallmatrix}P&Q\\0&P\end{smallmatrix}\right)$,
with $P\in\operatorname{GL}_3(\Z)$ and $Q\in\operatorname{Mat}_3(\Z)$.
\end{theorem}

\begin{theorem}\label{night:P1B.2}
\StatementRecord{night:P1B.2}{proved}{S6\_06\_BLIND: GPT-6 Astra (OpenAI, ChatGPT Pro mode, effort Pro), recorded re-derivation in the checkpoint dated 18 November 2026; also checked in this writing by \CurrentChecker}{Definitions and the displayed constructions}
A finite-order integral linear action on $\R^6$ admitting a
free finite affine torus lift has characteristic polynomial
among exactly thirty-eight cyclotomic products. Their linear
orders are
\[
 1,2,3,4,5,6,8,10,12.
\]
Exactly nineteen of these polynomials admit a commuting complex
structure on $\R^6$. These are a preliminary spectral list;
no assertion about global gluings is inferred from the list.
\end{theorem}

\begin{proposition}[the previously open period decision]\label{maps:src:F.3.3}
\StatementRecord{maps:src:F.3.3}{proved}{Claude Fable 5.2 (Anthropic), effort high, 16 November 2026; also checked in this writing by \CurrentChecker}{\cref{maps:src:F.3.4}, \cref{maps:src:F.3.7}, \cref{maps:src:F.3.8}, \cref{maps:src:F.3.9}, \cref{maps:src:F.3.10}, \cref{maps:src:F.3.11}}


Theorem \ref{maps:src:F.3.11} now gives a positive answer for every topological
survivor of F.2: there is a holomorphic period family on the
punctured sphere, with nondegenerate imaginary part in every
local normalized period chart and a bounded holomorphic
correction at its primitive maximal cusp.  The proof is universal
and does not require that the remaining integral classification
be finite.  The proposed polarized abelian-surface source is
excluded by Theorem \ref{maps:src:F.3.4}; the universal construction instead
uses the explicit unpolarized LSAP domain and the relative
section principle in \ref{maps:src:F.3.8}--\ref{maps:src:F.3.11}.  Simultaneous good reduction
at every finite puncture is a stronger compactification condition
than the literal F.3 statement and is not inferred from this
argument.  It is already proved for the explicit compact
fourfold of F.4.

\end{proposition}

\begin{theorem}[exclusion of the abelian-surface source]\label{maps:src:F.3.4}
\StatementRecord{maps:src:F.3.4}{proved}{GPT-6 Astra (OpenAI, ChatGPT Pro mode), effort Pro, 2026-11-17; Claude Fable 5.2 (Anthropic), effort high, 2026-11-16; also checked in this writing by \CurrentChecker}{the hypotheses, notation and definitions stated here; \cite[Definition 4.2.4 and Proposition 4.2.5, p.~44; Theorem 4.2.6, p.~45]{ffDeligne}; \cref{maps:src:F.3.1}}


The Engel construction proposed in F.3, starting with a polarized
algebraic family of abelian surfaces with a toric-rank-two cusp,
taking a fibrewise degree-zero line bundle, and quotienting its
punctured bundle by a lifted translation, cannot have $b_1=b_2=0$
with the stipulated rank-three square-zero Kuhn cusp.
The two sections need not be proportional.

\end{theorem}

\begin{theorem}[the unpolarized period domain has LSAP]\label{maps:src:F.3.8}
\StatementRecord{maps:src:F.3.8}{proved}{GPT-6 Astra (OpenAI, ChatGPT Pro mode), effort Pro, 2026-11-17; Claude Fable 5.2 (Anthropic), effort high, 16 November 2026; also checked in this writing by \CurrentChecker}{the hypotheses, notation and definitions stated here; \cite[Definition 2.1, p.~4]{IVVOkaOneV5}; \cite[Proposition 7.13 and preceding paragraph, pp.~23--24]{IVVOkaOneV5}; \cref{maps:eq:F.O3}; \cref{maps:eq:F.O4}}


Let $D_3$ be the component of
\[
 \{F\in\operatorname{Gr}(3,\mathbb C^6):F\cap\overline F=0\}
\]
containing the standard complex structure, so that
$D_3\cong GL^+(6,\mathbb R)/GL(3,\mathbb C)$.
At every point, $D_3$
is dominated by a spanning tube of nine complex lines.  In particular,
it has the local spray approximation property (LSAP) and is Oka--1.
These assertions concern the unpolarized domain.

\end{theorem}

\begin{theorem}[holomorphic periods on every open base]\label{maps:src:F.3.9}
\StatementRecord{maps:src:F.3.9}{proved}{GPT-6 Astra (OpenAI, ChatGPT Pro mode), effort Pro, 2026-11-17; Claude Fable 5.2 (Anthropic), effort high, 16 November 2026; also checked in this writing by \CurrentChecker}{the hypotheses, notation and definitions stated here; \cref{maps:src:F.3.8}; \cite[Proposition 7.13(d), p.~24]{IVVOkaOneV5}}


Let $B$ be an open connected Riemann surface and
$\rho:\pi_1(B)\to GL^+(6,\mathbb R)$ any representation.
Its associated flat $D_3$-bundle admits a holomorphic section.
More generally, every continuous section which is holomorphic
near a compact Runge set can be deformed through sections to a
holomorphic section, with arbitrary uniform approximation there
and any prescribed finite-order interpolation at finitely many
points of that set.  For an integral representation, this gives
a holomorphic family of complex three-tori with the prescribed
lattice local system and a zero section over $B$.

\end{theorem}

\begin{lemma}[retaining isolated extension points]\label{maps:src:F.3.10}
\StatementRecord{maps:src:F.3.10}{proved}{GPT-6 Astra (OpenAI, ChatGPT Pro mode), effort Pro, 2026-11-17; Claude Fable 5.2 (Anthropic), effort high, 16 November 2026; also checked in this writing by \CurrentChecker}{the hypotheses, notation and definitions stated here; \cref{maps:eq:F.O5}; \cite[Theorem 1.2, preprint p.~3]{ffForstnericPrezelj02}; \cite[Theorem 2.13, p.~9]{IVVOkaOneV5}}


Let $R$ be an open Riemann surface, $V\to R$ a holomorphic vector
bundle, and $\mathscr U\subset\operatorname{Gr}_r(V)$ an open
subset whose projection onto $R$ is surjective.  Let $S\subset R$
be finite.  Suppose that over $R\setminus S$ this projection is
a holomorphic fibre bundle with an LSAP fibre.  Let $K\subset R$
be a compact Runge set containing $S$ in its interior.
Every continuous section of $\mathscr U$ which is holomorphic
near $K$ can be deformed through sections to a holomorphic
section, with arbitrary uniform approximation on $K$ and
interpolation of any given finite jets at finitely many points
of its interior.  No fibre-bundle or Serre-fibration assertion
is required at the points of $S$.

\end{lemma}

\begin{theorem}[the bounded-cusp period decision]\label{maps:src:F.3.11}
\StatementRecord{maps:src:F.3.11}{proved}{GPT-6 Astra (OpenAI, ChatGPT Pro mode), effort Pro, 2026-11-17; Claude Fable 5.2 (Anthropic), effort high, 16 November 2026; also checked in this writing by \CurrentChecker}{the hypotheses, notation and definitions stated here; \cref{maps:src:F.2.1}; \cref{maps:eq:F.O7}; \cref{maps:src:F.3.8}; \cref{maps:src:F.3.10}; \cref{maps:src:F.2.7}; \cref{maps:src:F.2.15}}


Let
$B^*=\mathbb P^1\setminus\{p_0,p_1,\ldots,p_k\}$,
where $k\ge1$, and let
$\rho:\pi_1(B^*)\to GL^+(6,\mathbb R)$ be a representation.
Assume that the monodromy $M_0$ at $p_0$ satisfies
$(M_0-I)^2=0$ and $\operatorname{rank}(M_0-I)=3$.
Then there is a holomorphic period family with this monodromy
and nondegenerate real transversality everywhere on $B^*$,
equivalently a nondegenerate imaginary part in each local
normalized period chart.
In a marking with $M_0=\left(\begin{smallmatrix}I&I\\0&I\end{smallmatrix}\right)$
its cusp period is
\[
 Z(t)=\frac{\log t}{2\pi i}I_3+C(t),
 \qquad C(t)\text{ holomorphic at }t=0.
 \tag{F.O6}\label{maps:eq:F.O6}
\]
Any prescribed finite jet of $C$ at zero may be imposed.
For an integral primitive cusp the marking can be integral.
Consequently every topological survivor of F.2 admits the
period family requested literally in F.3.  No good-reduction
condition at the other, still punctured, points is included
in this assertion.

\end{theorem}